\documentclass[11pt,a4paper]{article}
\usepackage[margin=2.2cm]{geometry}
\usepackage[utf8]{inputenc}
\usepackage[english]{babel}
\usepackage{amsmath,amssymb,mathtools}
\usepackage{graphicx}
\usepackage{booktabs}
\usepackage{longtable}
\usepackage{array}
\usepackage{hyperref}
\usepackage{placeins}
\usepackage{xcolor}
\usepackage{caption}
\usepackage{subcaption}
\usepackage{enumitem}
\usepackage{algorithm}
\usepackage{algpseudocode}
\usepackage{listings}
\usepackage{siunitx}
\usepackage{microtype}
\usepackage{multirow}
\hypersetup{colorlinks=true,linkcolor=blue,urlcolor=blue,citecolor=blue}

\lstdefinestyle{cfg}{
  basicstyle=\ttfamily\footnotesize,
  breaklines=true,
  columns=fullflexible,
  frame=single,
  rulecolor=\color{black!30},
  showstringspaces=false,
}
\lstset{style=cfg}

\algrenewcommand\algorithmicrequire{\textbf{Input:}}
\algrenewcommand\algorithmicensure{\textbf{Output:}}

\newcommand{\code}[1]{\texttt{\small #1}}
\newcommand{\Vm}{V_{\text{m}}}
\newcommand{\Iion}{I_{\mathrm{ion}}}
\newcommand{\Dten}{\mathbf{D}}
\newcommand{\grad}{\nabla}
\newcommand{\divg}{\nabla\!\cdot}
\newcommand{\Mmass}{\mathbf{M}}
\newcommand{\Kstif}{\mathbf{K}}
\newcommand{\Lop}{\mathbf{L}}
\newcommand{\Aimp}{\mathbf{A}}
\newcommand{\Bimp}{\mathbf{B}}
\newcommand{\Fres}{\mathsf{F}}
\newcommand{\R}{\mathbb{R}}
\newcommand{\norm}[1]{\left\lVert #1 \right\rVert}
\newcommand{\Vmsub}[1]{V_{\text{m},#1}}
\newcommand{\Iionsub}[1]{I_{\mathrm{ion},#1}}
\newcommand{\Iext}{I_{\mathrm{ext}}}
\newcommand{\pyconst}[1]{\textcolor{black!60!blue}{\texttt{\footnotesize #1}}}
\newcommand{\dictkey}[1]{\textcolor{black!60!green}{\texttt{\footnotesize #1}}}
\title{Solver specification:\\
       \texorpdfstring{\code{highOrderElectroActivationFoamImplicitPETScDistributed}}{highOrderElectroActivationFoamImplicitPETScDistributed}\\
       \large A high-order, implicit, Jacobian-free Newton--Krylov monodomain
       solver in OpenFOAM with a switchable Eigen/PETSc linear-algebra
       back-end and distributed-memory linear algebra}
\author{Pablo}
\date{\today}

\begin{document}
\maketitle

\begin{abstract}
This report specifies the algorithms implemented in the OpenFOAM solver:

\begin{center}
\code{highOrderElectroActivationFoamImplicitPETSc}.
\end{center}

The solver advances the monodomain equations
of cardiac electrophysiology coupled to the
ten Tusscher--Noble--Noble--Panfilov (TNNP, 2004) cellular ionic model. The solver
has several features:
\begin{itemize}
\item The diffusion operator is
discretised with the OpenFOAM cell-centred finite-volume framework,
optionally enhanced to a high-order Local
Regression Estimators (LRE) reconstruction;
\item The ionic source is evaluated either at cell
centres or at the cell quadrature points (required for high order);
\item There are different ways to integrate the reaction ODE equations:
\begin{itemize}
\item[(i)] at each quadrature point used to evaluate the ionic source;
\item[(ii)] only at the centres, reconstructing them afterwards wherever the ionic source has to be evaluated;
\item[(iii)] a front-aware hybrid that does (i) only on the propagating front and (ii) everywhere else (Section~\ref{sec:stateint});
\end{itemize}
\item Time integration follows a $\theta$ scheme (backward Euler or Crank--Nicolson);
\item the per-cell ionic
ODE uses the OpenFOAM adaptive RKF45 integrator with configurable
tolerances;
\item The mass operator can be \code{lumped} or high-order \code{consistent};
\item the resulting coupled nonlinear system can be solved by Picard fixed-point
iteration, a diagonal linearisation of $\Iion$, or a Jacobian-Free
Newton--Krylov (JFNK) method with Armijo backtracking line search,
an initial guess by linear extrapolation and a physical bracket \emph{clamp}
on the input to the ionic ODE;
\item the linear system of the assembled branches (Picard / diagonalIion) and the
JFNK inner system are solved with a \emph{switchable} back-end via the
\code{linearSolverBackend} key: either Eigen's sparse factorisations
(\code{SparseLU}/\code{BiCGSTAB}) and an in-house restarted GMRES with modified
Gram--Schmidt, or PETSc's KSP/PC solvers (including a GMRES over a matrix-free
\code{MatShell} for JFNK). Matrix assembly and storage always use
\code{Eigen::SparseMatrix};
\item activation times are recorded by linear interpolation
across the first threshold crossing and the benchmark points are sampled by
inverse-distance weighting.
\end{itemize}
 
All algorithmic parameters are exposed
through the \code{constant/spatialIntegrationProperties} and
\code{constant/timeIntegrationProperties} dictionaries, and the Python driver
\code{run\_convergence.py} generates these dictionaries from a single
configuration file. This solver integrates both linear-algebra back-ends (Eigen and
PETSc) in a single executable, selectable at run time with
\code{linearSolverBackend} / \code{jfnkLinearSolverBackend}
(Section~\ref{sec:linalg}). There is also a simpler sibling solver,
\code{highOrderElectroActivationFoamImplicit}, with an Eigen-only back-end and a
semi-implicit single-Picard-iteration scheme (Section~\ref{sec:optim} for the
performance differences).

This \code{...Distributed} variant additionally makes the linear algebra
\emph{genuinely distributed}: the mesh is decomposed with \code{decomposePar},
the operators are assembled with local rows and global columns and become PETSc
\code{MPIAIJ} matrices on \code{PETSC\_COMM\_WORLD}, and every matrix--vector
product goes through \code{MatMult} instead of Eigen. The algorithms of
Sections~\ref{sec:stateint}--\ref{sec:linalg} are unchanged;
Section~\ref{sec:distributed} describes what the decomposition requires, how a
parallel operator is verified, and the one inconsistency that remains --- a
stencil selection in the high-order library that is not invariant under
partitioning (Section~\ref{sec:dist-stencils}).
\end{abstract}

\tableofcontents
\bigskip
% ================================================================== %
\section{Notation}
% ================================================================== %

\begin{center}
\begin{tabular}{@{}lll@{}}
\toprule
Symbol & Meaning & Units \\
\midrule
$\Omega \subset \R^{d}$ & Spatial domain ($d \in \{2,3\}$) & \si{\meter} \\
$\partial\Omega$ & Insulated boundary (zero-flux Neumann) & --- \\
$t \in (0, T]$ & Time & \si{\second} \\
$\Vm(\mathbf{x}, t)$ & Transmembrane voltage & \si{\volt} \\
$\mathbf{s}(\mathbf{x}, t) \in \R^{17}$ & TNNP ionic states & mixed \\
$\Iion(\Vm, \mathbf{s})$ & Total ionic current density & \si{\volt\per\second} \\
$\Iext(\mathbf{x}, t)$ & External stimulus current density & \si{\ampere\per\meter\cubed} \\
$\chi$ & Membrane surface-to-volume ratio & \si{\per\meter} \\
$C_{m}$ & Membrane capacitance per unit area & \si{\farad\per\meter\squared} \\
$\Dten(\mathbf{x})$ & Bulk conductivity tensor & \si{\siemens\per\meter} \\
$\theta \in \{0.5, 1\}$ & Time-stepping parameter (CN or BE) & --- \\
$\Delta t$ & Outer time-step length & \si{\second} \\
$N_{c}$ & Number of mesh cells & --- \\
$\mathsf{V}^{n} \in \R^{N_{c}}$ & Vector of cell-averaged $\Vm$ at time $t^{n}$ & \si{\volt} \\
$\Mmass, \Kstif, \Lop$ & Mass, stiffness, scaled diffusion matrix & --- \\
$\Aimp, \Bimp$ & Implicit LHS / RHS operators of the $\theta$-scheme & --- \\
$\Fres(\mathsf{x})$ & Discrete nonlinear residual to be zeroed & --- \\
\bottomrule
\end{tabular}
\end{center}

% ================================================================== %
\section{Governing equations}
% ================================================================== %

\subsection{Monodomain PDE}

The bulk voltage is propagated by the standard monodomain
reaction--diffusion equation~\cite{ttp2004,niederer2012}. In physical
current-density notation one may write
\begin{equation}
  \chi \, C_{m} \, \frac{\partial \Vm}{\partial t}
   \;=\;
   \divg\!\bigl( \Dten \grad \Vm \bigr)
   - \chi C_m\,\Iion(\Vm, \mathbf{s})
   + \Iext
   \quad \text{in } \Omega \times (0, T],
   \label{eq:monodomain}
\end{equation}
subject to insulated boundaries
$\mathbf{n} \cdot ( \Dten \grad \Vm ) = 0$ on
$\partial \Omega \times (0, T]$, and initial condition
$\Vm(\mathbf{x}, 0) = -\SI{84}{\milli\volt}$. The conductivity tensor
$\Dten$ is symmetric positive-definite and either uniform or read as a
\code{volTensorField} from disk. The factor $\chi C_m$ multiplying
$\Iion$ in~\eqref{eq:monodomain} reflects the convention used by the
ionic-model interface: the solver stores \code{Iion} as a voltage rate
with dimensions \si{\volt\per\second}, not as an
\si{\ampere\per\meter\cubed} volumetric current.

Dividing~\eqref{eq:monodomain} by $\chi C_{m}$ gives the form actually
used inside the solver,
\begin{equation}
  \frac{\partial \Vm}{\partial t}
   \;=\;
   \mathcal{L} \Vm
   \;+\;
   s(\Vm, \mathbf{s}),
   \qquad
   \mathcal{L} \;=\; \frac{1}{\chi C_{m}}\,\divg(\Dten \grad \cdot),
   \quad
   s \;=\; -\Iion + \frac{\Iext}{\chi C_{m}}.
   \label{eq:semi-discrete}
\end{equation}
This is exactly what \code{rhsFromIion} assembles in the solver.

\subsection{Ionic ODE}

At every spatial sampling point the TNNP state vector
$\mathbf{s} \in \R^{17}$ evolves according to
\begin{equation}
  \frac{\mathrm{d}\mathbf{s}}{\mathrm{d}t}
   \;=\;
   \mathbf{f}_{\text{TNNP}}\!\bigl(\Vm, \mathbf{s}; \texttt{tissue}\bigr),
   \qquad \mathbf{s}(0) = \mathbf{s}_{0},
   \label{eq:tnnp}
\end{equation}
where \code{tissue} selects \code{epicardialCells}, \code{mCells} or
\code{endocardialCells}. The 17 tracked states are
\[
  \mathbf{s} = (V,\ K_{i},\ Na_{i},\ Ca_{i},\ Xr_{1},\ Xr_{2},\ Xs,\ m,\ h,\ j,\ d,\ f,\ fCa,\ s,\ r,\ g,\ Ca_{SR}).
\]
The total ionic current returned to the PDE is the algebraic
\code{Iion\_cm} of TNNP~\cite{ttp2004}.

\subsection{Manufactured MMS equation used for convergence transfer}

The convergence driver in
\code{tutorials\_electro/highOrderManufacturedFDAImplicitJfNK} solves a
manufactured problem with the same numerical structure as the
electro-activation solver: implicit monodomain-like PDE, local state
ODEs, nonlinear ionic source, optional LRE reconstruction for $\Vm$ and
$\Iion$, and selectable \code{Picard}/\code{diagonalIion}/\code{JFNK}
nonlinear coupling. The analytical fields are
\begin{align}
  \Vm^{\star}(\mathbf{x},t) &= \sqrt{1+t}\,F(\mathbf{x}),\\
  u_1^{\star}(\mathbf{x},t) &= (1+t)\,G(\mathbf{x})+\Vm^{\star}(\mathbf{x},t),\\
  u_2^{\star}(\mathbf{x},t) &=
      \frac{1}{(1+t)\sqrt{G(\mathbf{x})}},\\
  u_3^{\star}(\mathbf{x},t) &= 0,
\end{align}
with, for example,
$F=\cos(\pi x)\cos(2\pi y)$ and $G=1+x\,y^2$ in 2D. In 3D,
$F=\cos(\pi x)\cos(2\pi y)\cos(3\pi z)$ and
$G=1+x\,y^2 z^3$. The state equations are
\begin{align}
  \dot u_1 &=
      (u_1+u_3-\Vm)^2 u_2^2
    + \frac{1}{2}(u_1+u_3-\Vm)u_2^2(\Vm-u_3),\\
  \dot u_2 &= -(u_1+u_3-\Vm)u_2^3,\\
  \dot u_3 &= 0.
\end{align}
The manufactured ionic source is
\begin{equation}
  \Iion^{\mathrm{MMS}}(\Vm,\mathbf{u})
  =
  -\frac{1}{2}(u_1+u_3-\Vm)u_2^2(\Vm-u_3)
  + \frac{\beta}{\chi C_m}(\Vm-u_3),
\end{equation}
and the PDE source used by the manufactured solver is
$-\Iion^{\mathrm{MMS}}$. The coefficient $\beta$ is the spatial
eigenvalue of the trigonometric manufactured field under the constant
conductivity tensor:
\[
  \beta =
  \begin{cases}
    -\pi^2 D_{xx}, & d=1,\\
    -\pi^2(D_{xx}+4D_{yy}), & d=2,\\
    -\pi^2(D_{xx}+4D_{yy}+9D_{zz}), & d=3.
  \end{cases}
\]
Thus the MMS case is not intended to reproduce TNNP physiology; it is a
controlled problem with a known exact solution that exercises the same
discrete operators and nonlinear solution paths used by this solver.

\subsection{External stimulus}

The external current is a time- and region-windowed step:
\begin{equation}
  \Iext(\mathbf{x}, t) \;=\;
  \begin{cases}
    I_{0} & \mathbf{x} \in \Omega_{\text{stim}}^{(b)},
            \;\; t \in [t_{0}^{(b)},\, t_{0}^{(b)} + \tau],
            \;\; b = 1, \dots, B, \\
    0     & \text{otherwise},
  \end{cases}
\end{equation}
with $I_{0} = \SI{50000}{\ampere\per\meter\cubed}$ and
$\tau = \SI{2}{\milli\second}$ as Niederer-benchmark defaults.

% ================================================================== %
\section{Time discretisation: the $\theta$-scheme}
% ================================================================== %

Time is advanced from $t^{n}$ to $t^{n+1} = t^{n} + \Delta t$ by
\begin{equation}
  \frac{\mathsf{V}^{n+1} - \mathsf{V}^{n}}{\Delta t}
   \;=\;
   \theta\, \Lop\, \mathsf{V}^{n+1}
   + (1-\theta)\, \Lop\, \mathsf{V}^{n}
   + \theta\, \mathbf{s}^{n+1}
   + (1-\theta)\, \mathbf{s}^{n}.
   \label{eq:theta}
\end{equation}

\begin{center}
\begin{tabular}{@{}lcl@{}}
\toprule
\code{implicitScheme} & $\theta$ & Scheme \\
\midrule
\code{backwardEuler} & $1$   & First-order, L-stable \\
\code{crankNicolson} & $1/2$ & Second-order, A-stable \\
\bottomrule
\end{tabular}
\end{center}

After multiplication by $\Delta t^{-1}$ the implementation forms
\begin{equation}
  \bigl[\Mmass/\Delta t \,-\, \theta\, \Lop\bigr]\, \mathsf{V}^{n+1}
   \;=\;
   \bigl[\Mmass/\Delta t \,+\, (1-\theta)\, \Lop\bigr]\, \mathsf{V}^{n}
   \;+\; \theta\, \mathbf{s}^{n+1}
   \;+\; (1-\theta)\, \mathbf{s}^{n},
   \label{eq:implicit-system}
\end{equation}
with the two Eigen \code{SparseMatrix<scalar,RowMajor>} operators
\begin{equation}
  \Aimp \;=\; \Mmass/\Delta t \,-\, \theta\, \Lop,
   \qquad
  \Bimp \;=\; \Mmass/\Delta t \,+\, (1-\theta)\, \Lop.
   \label{eq:AB}
\end{equation}
$\Aimp$ and $\Bimp$ are rebuilt at every time step by copying $\Mmass$,
scaling by $\Delta t^{-1}$ and adding $\pm \theta \Lop$. When
$\theta = 1$ (BE) the explicit half of the diffusion vanishes and
$\Bimp = \Mmass/\Delta t$. The scaled diffusion is
$\Lop = \Kstif / (\chi C_{m})$.

% ================================================================== %
\section{Spatial discretisation}
% ================================================================== %

\subsection{Standard orthogonal FV Laplacian}

\code{assembleStandardOrthogonalStiffnessMatrix} computes the classical
5/7-point Laplacian. Per internal face $f$ with face area $\mathbf{S}_{f}$,
\begin{equation}
  a_{f} \;=\;
   \frac{ \lvert \mathbf{S}_{f} \rvert\,
          \bigl( \mathbf{n}_{f} \cdot \Dten_{f} \cdot \hat{\mathbf{d}}_{f} \bigr) }
        { \lvert \mathbf{d}_{f} \rvert },
   \quad
   \Dten_{f} = \tfrac{1}{2}\,(\Dten_{\text{own}} + \Dten_{\text{nei}}),
\end{equation}
contributing $-a_{f}/V_{\text{own}}$ to $\Kstif[\text{own},\text{own}]$,
$+a_{f}/V_{\text{own}}$ to $\Kstif[\text{own},\text{nei}]$, and the
transposed pair to the neighbour row. \code{zeroGradient}/\code{empty}
patches contribute zero flux. \code{fixedValue} patches contribute to
the owner diagonal only.

\subsection{High-order LRE diffusion}

When \code{useHighOrder\_Vm = true} and $d>1$,
\code{assembleHighOrderStiffnessMatrix} integrates the flux on a face
quadrature with reconstructed gradients:
\begin{equation}
  \Phi_{f} \;=\;
   \sum_{q \in \mathcal{Q}_{f}}
       w_{q}\, \lvert \mathbf{S}_{f} \rvert\,
       \mathbf{n}_{f} \cdot \Dten_{f} \cdot
       \Bigl(
         \sum_{j \in \mathcal{S}_{f}}
            \mathbf{g}_{f,q,j}\, \mathsf{V}_{j}
       \Bigr),
   \label{eq:HO-flux}
\end{equation}
where $\mathcal{S}_{f}$ is the LRE face stencil ($\sim 15$--$50$ cells)
and $\mathbf{g}_{f,q,j} = $~\code{QRGradFaceGPCoeffs}. The
contribution of each $\mathsf{V}_{j}$ is added to row \code{own} with
weight $+\Phi_{f,j}/V_{\text{own}}$ and to row \code{nei} with weight
$-\Phi_{f,j}/V_{\text{nei}}$, yielding a negative semi-definite
Laplacian.

\subsection{Mass matrices}

Two options:
\paragraph{Lumped (\code{massMatrix lumped}).}
\code{assembleDiagonalMassMatrix} writes $\Mmass = I$.

\paragraph{Consistent HO (\code{massMatrix consistent}).}
\code{assembleConsistentMassMatrixHO} fills
\begin{equation}
  \Mmass[i, j] \;=\;
   \sum_{q \in \mathcal{Q}_{i}}
       \frac{w_{q}}{\sum_{q'} w_{q'}}\,
       \phi_{j}(\mathbf{x}_{q}),
   \quad
   \phi_{j}(\mathbf{x}_{q})
    = \delta_{ij}
    + \mathbf{g}_{i,j} \cdot \mathbf{d}_{q}
    + \tfrac{1}{2}\, \mathbf{H}_{i,j} \!:\! (\mathbf{d}_{q} \otimes \mathbf{d}_{q})
    + \tfrac{1}{6}\, \mathbf{T}^{(3)}_{i,j} \!\vdots\!
        (\mathbf{d}_{q} \otimes \mathbf{d}_{q} \otimes \mathbf{d}_{q}),
   \label{eq:HO-mass}
\end{equation}
with $\mathbf{d}_{q} = \mathbf{x}_{q} - \mathbf{C}_{i}$. A safety guard
silently downgrades to the lumped form when
\code{useHighOrder\_Vm = false} or $d \le 1$.

\subsection{High-order reconstruction at ionic quadrature points}

If \code{useHighOrder\_Iion = true}, the ionic ODE is integrated at
cell quadrature points. Vm is reconstructed there by the same Taylor
expansion as in~\eqref{eq:HO-mass}:
\begin{equation}
  \Vm(\mathbf{x}_{q})
   \;\approx\;
   \Vmsub{c} + (\grad \Vm)_{c} \cdot \mathbf{d}_{q}
   + \tfrac{1}{2}\, \mathbf{H}_{c} \!:\! (\mathbf{d}_{q} \otimes \mathbf{d}_{q})
   + \tfrac{1}{6}\, \mathbf{T}^{(3)}_{c} \!\vdots\!
       (\mathbf{d}_{q} \otimes \mathbf{d}_{q} \otimes \mathbf{d}_{q}),
   \label{eq:HO-vm}
\end{equation}
implemented in \code{reconstructVmAtIionIntegrationPoints}.

\subsection{High-order cell averaging of $\Iion$}

\code{averageIionIntegrationPointsToCells} contracts the per-QP $\Iion$
back to cells:
\begin{equation}
  \Iionsub{i} \;=\;
   \frac{\sum_{q \in \mathcal{Q}_{i}} w_{q}\, \Iionsub{q}}
        {\sum_{q \in \mathcal{Q}_{i}} w_{q}}.
\end{equation}

% ================================================================== %
\section{State integration at the ionic quadrature points}
\label{sec:stateint}
% ================================================================== %

The high-order quadrature~\eqref{eq:HO-vm} evaluates $\Iion(\Vm,\mathbf{s})$
at the cell quadrature points $\mathbf{x}_q$, so it needs the ionic state
vector $\mathbf{s}$ there as well. \emph{How} the states reach the quadrature
points is governed by \dictkey{stateIntegrationMode} and, orthogonally, by
\dictkey{frontAwareHybrid}. This choice turns out to be decisive for a
propagating cardiac front, and this section documents both the algorithm and
the reason it matters.

% ------------------------------------------------------------------ %
\subsection{Two modes}
% ------------------------------------------------------------------ %

Let $\mathcal{Q}_i=\{\mathbf{x}_q\}$ be the quadrature set of cell $\Omega_i$
with weights $w_q$, and let $\mathbf{s}_c$ denote a cell-centred state.

\paragraph{\code{gaussPointODE} (legacy).}
The stiff ODE~\eqref{eq:tnnp} is integrated \emph{independently at every
quadrature point}, driven by the locally reconstructed potential:
\begin{equation}
  \mathbf{s}_q^{\,n+1}
   = \mathrm{ODEsolve}\!\bigl(\mathbf{s}_q^{\,n},\,\Vm(\mathbf{x}_q);\Delta t\bigr),
  \qquad
  \Iionsub{q}=\phi\bigl(\Vm(\mathbf{x}_q),\mathbf{s}_q^{\,n+1}\bigr).
  \label{eq:gpode}
\end{equation}
Cost: $N_q$ stiff integrations per cell.

\paragraph{\code{cellCentredReconstruct}.}
The ODE is integrated \emph{once per cell} at the cell centre, and the
resulting cell-centred states are reconstructed to the quadrature points by a
dedicated LRE (\code{LREInterp\_states}, order set by
\dictkey{highOrderCoeffs/LRECoeffs\_states}):
\begin{equation}
  \mathbf{s}_c^{\,n+1}
   = \mathrm{ODEsolve}\!\bigl(\mathbf{s}_c^{\,n},\,\overline{V}_{m,c};\Delta t\bigr),
  \quad
  \widehat{\mathbf{s}}(\mathbf{x}_q)
   = \mathcal{R}^{(p)}\!\bigl[\{\mathbf{s}_{c'}^{\,n+1}\}\bigr](\mathbf{x}_q),
  \quad
  \Iionsub{q}=\phi\bigl(\Vm(\mathbf{x}_q),\widehat{\mathbf{s}}(\mathbf{x}_q)\bigr).
  \label{eq:ccrec}
\end{equation}
Cost: one stiff integration per cell -- a factor $N_q$ cheaper, which for the
expensive TNNP integrations is the motivation. The two modes differ only in the
argument of the non-linear map $\Vm\!\mapsto\!\mathbf{s}$: pointwise
$\Vm(\mathbf{x}_q)$ in~\eqref{eq:gpode}, versus the cell average
$\overline{V}_{m,c}$ in~\eqref{eq:ccrec}.

% ------------------------------------------------------------------ %
\subsection{Front-aware hybrid}
\label{sec:hybrid}
% ------------------------------------------------------------------ %

\dictkey{frontAwareHybrid} (default \code{false}, orthogonal to
\dictkey{stateIntegrationMode}) combines the two modes: it solves the
ODE at the Gauss points in the thin band of \emph{front} cells and at the cell
centre (with reconstruction) everywhere else. A cell $i$ is a \emph{front} cell
when its sub-cell $\Vm$ range exceeds \dictkey{hybridFrontVmRange} or it is
stimulated,
\begin{equation}
  \text{front}_i \;\Longleftrightarrow\;
  \Bigl(\max_q\Vm(\mathbf{x}_q)-\min_q\Vm(\mathbf{x}_q)\Bigr)
     > \dictkey{hybridFrontVmRange}
  \;\;\vee\;\; (\Iext)_i\neq 0 .
  \label{eq:frontflag}
\end{equation}
The persistent states are carried at the Gauss points (as in
\code{gaussPointODE}); the smooth-cell reconstruction derives its cell values
by quadrature-averaging that carrier. Per nonlinear evaluation the branch is:
\begin{enumerate}[topsep=2pt,itemsep=1pt,label=(\alph*)]
  \item front cells: one ODE per Gauss point, local $\Vm$ (as~\eqref{eq:gpode});
  \item smooth cells: one ODE at the cell centre from the cell-average state;
  \item reconstruct the smooth-cell states to their Gauss points, front Gauss
        states left untouched;
  \item evaluate $\Iion$ at all Gauss points and average to the cell.
\end{enumerate}

\paragraph{Freezing per time step (avoiding a limit cycle).}
The front classification~\eqref{eq:frontflag} is a \emph{solution-dependent
discrete switch}. If recomputed inside the nonlinear loop, cells near the
threshold flip front$\leftrightarrow$smooth between iterations; because the two
branches compute $\Iion$ by different paths, the source term jumps discretely
and the Picard fixed-point map becomes discontinuous, producing a non-convergent
\emph{period-2 limit cycle} (observed as $\text{iter}_{k+2}\equiv\text{iter}_k$
in \code{nonlinearResiduals.dat}). The classification is therefore computed
\emph{once per time step} from the committed $\Vm^{n}$ and \emph{frozen} for the
whole nonlinear solve, and the band is dilated by \dictkey{hybridFrontDilation}
cell layers (via \code{mesh.cellCells()}) so the front stays inside the
Gauss-resolved zone as it advances. Freezing is safe because the front moves far
less than one cell per step (for the 2D benchmark, $\sim\!\SI{0.09}{mm/ms}$, i.e.
$\sim\!0.009$ cell per step at $\Delta t=\SI{e-4}{s}$, $\Delta x=\SI{0.5}{mm}$).

% ------------------------------------------------------------------ %
\subsection{State reconstruction limiter}
\label{sec:limiter}
% ------------------------------------------------------------------ %

High-order ($p3$) reconstruction of the stiff TNNP states across an
under-resolved front overshoots (Gibbs), and a state driven far outside its
physical range can overflow the ionic algebra (observed as a PETSc
\code{SIGFPE}). \dictkey{stateReconstructLimiter} (default \code{true}, active
for \code{cellCentredReconstruct} and the hybrid smooth cells, independent of
\dictkey{frontAwareHybrid}) applies a Barth--Jespersen-style bound: each
reconstructed value is clipped to the range of the cell and its face neighbours,
\begin{equation}
  \widehat{s}(\mathbf{x}_q)
  \;\leftarrow\;
  \mathrm{clip}\!\Bigl(\widehat{s}(\mathbf{x}_q),\;
     \min_{c'\in\{i\}\cup\mathcal{N}(i)} s_{c'},\;
     \max_{c'\in\{i\}\cup\mathcal{N}(i)} s_{c'}\Bigr).
  \label{eq:limiter}
\end{equation}
No new extremum is created, so the reconstructed states stay within the
physically realised (ODE-evolved) range and the ionic algebra never sees a
pathological input; in smooth regions the reconstruction already lies within the
range and the limiter is inactive, preserving the order there.

% ------------------------------------------------------------------ %
\subsection{Dictionary keys}
% ------------------------------------------------------------------ %

\begin{center}
\begin{tabular}{@{}lll@{}}
\toprule
Key (in \code{implicitHOCoeffs}) & Meaning & Default \\
\midrule
\code{stateIntegrationMode} & \code{gaussPointODE} / \code{cellCentredReconstruct} & \code{cellCentredReconstruct} \\
\code{frontAwareHybrid} & Gauss ODEs on the front, reconstruct elsewhere & \code{false} \\
\code{hybridFrontVmRange} & sub-cell $\Vm$ range flagging a front cell [V] & $5\times10^{-3}$ \\
\code{hybridFrontDilation} & cell layers added around the front band & $2$ \\
\code{stateReconstructLimiter} & Barth--Jespersen limiter on the reconstruction & \code{true} \\
\bottomrule
\end{tabular}
\end{center}
\code{LRECoeffs\_states} (in \code{highOrderCoeffs}) sets the reconstruction
order; it is required by \code{cellCentredReconstruct} and by the hybrid.

% ================================================================== %
\section{Nonlinear solver}
% ================================================================== %

At each outer time step the unknown is the new voltage vector
$\mathsf{x}=\mathsf{V}^{n+1}$. For any candidate $\mathsf{x}$ the solver
first reconstructs the corresponding $\Vm$ field, then calls
\code{evaluateNonlinearFields}. That routine resets the ionic states to
their time-step-start values $\mathbf{s}^{n}$, integrates the local ODEs
from $t^n$ to $t^{n+1}$ using the candidate voltage, obtains
$\Iion(\mathsf{x})$, and optionally averages quadrature-point values
back to cells. This gives the nonlinear residual
\begin{equation}
  \Fres(\mathsf{x})
   \;=\;
   \Aimp\, \mathsf{x}
   \;-\;
   \mathsf{r}\!\bigl( \Iion(\mathsf{x}) \bigr),
   \qquad
   \mathsf{x} = \mathsf{V}^{n+1},
   \label{eq:Fdef}
\end{equation}
with
\begin{equation}
  \mathsf{r}(\Iion)
  =
  \Bimp\,\mathsf{V}^{n}
  + \theta\,\Bigl[-\Iion(\mathsf{x})+\Iext^{n+1}/(\chi C_m)\Bigr]
  + (1-\theta)\,\Bigl[-\Iion^{n}+\Iext^{n}/(\chi C_m)\Bigr].
  \label{eq:rhsIion}
\end{equation}
The three nonlinear methods differ only in how they drive
$\Fres(\mathsf{x})$ toward zero.

\subsection{Shared residuals and stopping tests}

Every method records the same diagnostics:
\begin{align}
  R_{\mathrm{cpl}} &=
    \frac{\|\Fres(\mathsf{x}_{k})\|_2}
         {\|\mathsf{x}_{k}\|_2+\epsilon},\\
  R_{\Vm} &=
    \frac{\|\mathsf{V}_{k}-\mathsf{V}_{k-1}\|_2}
         {\|\mathsf{V}_{k-1}\|_2+\epsilon},\\
  R_{\Iion} &=
    \frac{\|I_{\mathrm{ion},k}-I_{\mathrm{ion},k-1}\|_2}
         {\|I_{\mathrm{ion},k-1}\|_2+\epsilon},\\
  R_{\mathbf{s}} &=
    \max_j
    \frac{\|\mathbf{s}_{j,k}-\mathbf{s}_{j,k-1}\|_2}
         {\|\mathbf{s}_{j,k-1}\|_2+\epsilon}.
\end{align}
The actual implementation uses \code{relativeL2Difference} for
$R_{\Vm}$ and $R_{\Iion}$, and \code{maxStateRelativeL2Difference} for
the state maximum. A nonlinear iteration can stop only after
\code{minNonlinearIterations}; convergence then requires
\code{nonlinearTolerance}, \code{nonlinearVmTolerance} and
\code{nonlinearIionTolerance}. The state criterion
\code{nonlinearStatesTolerance} is included only when
\code{nonlinearRequireStatesConvergence=true}.

\subsection{Picard fixed-point method}

The Picard branch treats the source as lagged during the PDE solve. At
iteration $k$ the states and ionic source are computed from
$\mathsf{x}_k$, then the linear PDE
\begin{equation}
  \Aimp\,\mathsf{V}_{*}
  =
  \mathsf{r}\!\bigl(\Iion(\mathsf{x}_{k})\bigr)
\end{equation}
is solved with the Eigen sparse solver selected by \code{linearSolver}.
The next candidate is relaxed,
\begin{equation}
  \mathsf{x}_{k+1}
  =
  \mathsf{x}_{k}
  + \omega(\mathsf{V}_{*}-\mathsf{x}_{k}),
  \qquad
  \omega=\code{nonlinearRelaxation}.
\end{equation}
After the update the solver re-evaluates states and $\Iion$ at
$\mathsf{x}_{k+1}$ and checks the residuals above. In this
electro-activation solver, unlike a pure MMS Picard study, a failed
Picard solve rolls back unless \code{nonlinearAcceptUnconverged=true}.

\paragraph{Advantages.}
Picard is simple, cheap per nonlinear iteration, and uses the same
linear PDE solver as the classical implicit formulation. It is a useful
baseline and is often adequate when $\Delta t$ is small or the ionic
source is weakly nonlinear over one time step.

\paragraph{Disadvantages.}
The source derivative with respect to $\Vm$ is ignored in the PDE
matrix, so convergence can be slow near the action-potential upstroke or
for large $\Delta t$. A small residual may require many iterations, and
under-relaxation may be needed for robustness.

\subsection{Diagonal \texorpdfstring{$\Iion$}{Iion} linearisation}

The \code{diagonalIion} branch is Picard plus a local finite-difference
approximation of the source derivative. Define
\[
  S(\Vm) = -\Iion(\Vm,\mathbf{s}(\Vm)) + \Iext/(\chi C_m).
\]
The solver approximates, cell by cell,
\begin{equation}
  S'_i(V_{m,k})
  \;\approx\;
  -\frac{
      I_{\mathrm{ion},i}(V_{m,k}+\varepsilon_d)
      - I_{\mathrm{ion},i}(V_{m,k})
    }{\varepsilon_d},
  \qquad
  \varepsilon_d=\code{diagonalIionEpsilon}.
\end{equation}
The linearised residual
$\Aimp\mathsf{x}-\theta S(\mathsf{x})$ gives the system
\begin{equation}
  \left(\Aimp-\theta\,\mathrm{diag}(S'_k)\right)\mathsf{x}_{k+1}
  =
  \mathsf{r}\!\bigl(\Iion(\mathsf{x}_{k})\bigr)
  - \theta\,\mathrm{diag}(S'_k)\,\mathsf{x}_{k}.
  \label{eq:diagonalIion}
\end{equation}
The update can still be relaxed with \code{nonlinearRelaxation}, and
the same convergence and rollback logic as Picard is used.

\paragraph{Advantages.}
It captures the dominant local stiffness of $\Iion(\Vm)$ without
building a full Jacobian. It often reduces the number of nonlinear
iterations relative to Picard while preserving a sparse scalar PDE
matrix.

\paragraph{Disadvantages.}
Only the cell-local diagonal sensitivity is included. Couplings through
ODE states, high-order quadrature averaging and spatial reconstruction
are not represented in the matrix. The finite-difference step
\code{diagonalIionEpsilon} must be chosen above ODE/noise levels but
small enough to be a useful derivative.

\subsection{Jacobian-Free Newton--Krylov}

JFNK applies Newton to~\eqref{eq:Fdef}:
\[
  \mathbf{J}(\mathsf{x}_k)\delta_k
  =
  -\Fres(\mathsf{x}_k),
  \qquad
  \mathsf{x}_{k+1}=\mathsf{x}_k+\alpha_k\delta_k,
\]
but never assembles $\mathbf{J}$. GMRES sees only matrix-vector
products approximated by
\begin{equation}
  \mathbf{J}(\mathsf{x}_{k})\, \mathbf{v}
   \;\approx\;
   \frac{\Fres(\mathsf{x}_{k} + \varepsilon\, \mathbf{v})
         - \Fres(\mathsf{x}_{k})}{\varepsilon},
   \qquad
   \varepsilon \;=\; \varepsilon_{0}\,
                     \frac{1 + \norm{\mathsf{x}_{k}}_{2}}
                          {\norm{\mathbf{v}}_{2}},
   \label{eq:fdjac}
\end{equation}
with $\varepsilon_{0}=\code{jfnkEpsilon}$. Every evaluation of
$\Fres$ re-runs the ODEs and recomputes $\Iion$ from the same
time-step-start state snapshot. This is expensive but gives a genuinely
coupled nonlinear residual.

\subsection{JFNK inner Krylov}
\label{sec:jfnk-krylov}

The Newton linear system $\mathcal{J}_k\delta=-\Fres_k$ is
solved with a back-end chosen by \code{jfnkLinearSolverBackend}:
\begin{itemize}[nosep]
\item \textbf{Eigen} --- \code{solveGMRES} implements GMRES(m) with modified
  Gram--Schmidt, using at most \code{jfnkMaxKrylovIterations} basis vectors per
  restart and \code{jfnkMaxRestarts} restarts.
\item \textbf{PETSc} --- \code{KSPGMRES} over a \code{MatShell} whose
  matrix-vector product is the finite-difference wrapper; the KSP object is
  persistent across steps.
\end{itemize}
In both cases the stopping test is controlled by \code{jfnkLinearTolerance} and the
method is matrix-free: the linear operator is the finite-difference
shell of~\eqref{eq:fdjac}, not a stored sparse matrix.

\paragraph{Preconditioner (optional).}
With \code{jfnkPreconditioner=diffusion} a right preconditioner
$\mathbf{M}^{-1}$ is applied with $\mathbf{M}=\Aimp$ (the linear, diffusive, SPD part of the
Jacobian), factorised once as an incomplete ILU and reused; each Krylov step then
costs a cheap \emph{triangular solve} instead of an ODE
integration. The default \code{none} reproduces the unpreconditioned path.
See Section~\ref{sec:optim-A3}.

\paragraph{Advantages.}
JFNK sees the full nonlinear dependence
$\Vm\mapsto\mathbf{s}(\Vm)\mapsto\Iion(\Vm,\mathbf{s})$ and usually has much better
nonlinear convergence than Picard in stiff regimes. It also avoids
explicitly forming a dense or complicated Jacobian.

\paragraph{Disadvantages.}
Every Krylov vector requires an extra residual evaluation, and hence a
fresh ODE integration in every cell or quadrature point. The method is sensitive
to the finite-difference perturbation, to the ODE tolerances and to the
line-search choices. Without a preconditioner
(\code{jfnkPreconditioner=none}) the number of matvecs --- and therefore of
ODE integrations --- can be high; the diffusion preconditioner
(Section~\ref{sec:jfnk-krylov}) reduces it markedly. Empirically, for the
benchmark of this report JFNK turns out to be considerably more expensive than Picard even with
a preconditioner (Section~\ref{sec:exp}).
\subsection{Armijo backtracking line search}

The raw Newton update may increase the nonlinear residual when the
linear model is poor. If \code{jfnkLineSearch=true}, the solver tries
\[
  \mathsf{x}_{k+1}=\mathsf{x}_{k}+\alpha\delta_{k},
  \qquad
  \alpha=\code{nonlinearRelaxation},\;
  \alpha/2,\;\alpha/4,\dots
\]
until the Armijo sufficient-decrease condition holds:
\begin{equation}
  \norm{\Fres(\mathsf{x}_{k}+\alpha\delta_{k})}_2^2
  \le
  (1-2c\alpha)\,\norm{\Fres(\mathsf{x}_{k})}_2^2,
  \qquad c=\code{jfnkArmijoC}.
  \label{eq:armijo}
\end{equation}
Backtracking stops after \code{jfnkLineSearchMaxIter} attempts or when
the step length reaches \code{jfnkLineSearchAlphaMin}. This globalises Newton:
large Newton steps are accepted when they reduce the residual, otherwise
they are damped.

\subsection{Initial guess, ODE clamp and rollback}

The first guess for the time step is either $\mathsf{V}^{n}$ or the
linear extrapolation
\begin{equation}
  \mathsf{x}_{0}
  =
  \mathsf{V}^{n}
  + \frac{\Delta t}{\Delta t_{\mathrm{prev}}}
    \left(\mathsf{V}^{n}-\mathsf{V}^{n-1}\right),
  \label{eq:extrap}
\end{equation}
selected by \code{jfnkInitGuessOrder}. It is clamped to
$[\code{jfnkInitGuessVmMin},\code{jfnkInitGuessVmMax}]$. The same
bracket can also be applied to voltage values passed into the ODE solver
through \code{jfnkClampODEInput}; this protects the ionic model during
JFNK finite-difference probes.

If the nonlinear loop reaches its iteration budget without satisfying
the convergence tests, \code{nonlinearAcceptUnconverged} decides whether
the last iterate is committed. With the default \code{false}, the solver
rolls back $\Vm$, $\Iion$ and all ionic states to the beginning of the
time step. Rollback prevents a bad nonlinear solve from contaminating
the remaining activation simulation.

% ================================================================== %
\section{Linear solver: Eigen and PETSc back-ends}
\label{sec:linalg}
% ================================================================== %

Matrix assembly and storage ($\Mmass$, $\Kstif$, $\Lop$,
$\Aimp$) \emph{always} use \code{Eigen::SparseMatrix}. Solving the linear
system of the assembled branches (Picard / diagonalIion) is handled by the
\code{solveSparseSystem} wrapper, with a back-end chosen by \code{linearSolverBackend}:

\begin{center}
\begin{tabular}{@{}lll@{}}
\toprule
\code{linearSolverBackend} & \code{linearSolver} & Implementation \\
\midrule
\code{Eigen} & \code{SparseLU}  & \code{Eigen::SparseLU} (direct LU) \\
\code{Eigen} & \code{BiCGSTAB}  & \code{Eigen::BiCGSTAB}+\code{IncompleteLUT} \\
\code{PETSc} & (mapped)         & \code{KSP}+\code{PC} (e.g.\ \code{gmres}+\code{ilu},
   \code{preonly}+\code{lu}, \code{cg}+\code{gamg}) \\
\bottomrule
\end{tabular}
\end{center}

With the PETSc back-end, the Eigen matrix is copied to a PETSc \code{Mat} and the
KSP/PC type is taken from \code{petscLinearKspType} / \code{petscLinearPcType}
(the latter accepts \code{ilu}, \code{lu}, \code{jacobi}, \code{gamg}, \code{hypre},
\dots). All of them share \code{tolerance} and \code{maxIterations}.

\paragraph{Factorisation reuse.}
Since $\Aimp=\Mmass/\Delta t-\theta\Lop$ is constant as long as $\Delta t$ does not change,
the matrix is assembled once and the solver (PETSc KSP or Eigen
factorisation) is kept for the whole run, refreshed only when needed according to
an update policy \{Rebuild, ValuesOnly, Reuse\}. This avoids
re-factorising at every nonlinear iteration (Section~\ref{sec:optim-B}).

\paragraph{JFNK.}
The matrix-free JFNK branch does not instantiate these assembled wrappers; it uses
its own Krylov solver (Section~\ref{sec:jfnk-krylov}) over the finite-difference
Jacobian shell, with back-end \code{jfnkLinearSolverBackend} (Eigen or PETSc).

% ================================================================== %
\section{Ionic ODE integration}
% ================================================================== %

For each cell or cell QP, the TNNP ODE~\eqref{eq:tnnp} is advanced by
the OpenFOAM adaptive integrator selected via the \code{solver} key.
Default \code{RKF45}~\cite{hairer2}. Keys read from
\code{constant/timeIntegrationProperties}:

\begin{center}
\begin{tabular}{@{}lll@{}}
\toprule
Key & Meaning & Default \\
\midrule
\code{solver} & Concrete ODE integrator & \code{RKF45} \\
\code{absTol} & Absolute tolerance per state & $10^{-9}$ \\
\code{relTol} & Relative tolerance per state & $10^{-6}$ \\
\code{maxSteps} & Cap on sub-steps per outer call & $10^{4}$ \\
\code{initialODEStep} & Initial sub-step (ms) & $10^{-5}$ \\
\bottomrule
\end{tabular}
\end{center}

\code{relTol} drives the noise floor of the matrix-free Jacobian:
$\eta_{\Fres} \sim $~\code{relTol}. Tightening \code{relTol} buys
a lower attainable outer tolerance at the cost of more RKF45 sub-steps
per Newton iteration.

% ================================================================== %
\section{TNNP numerical protection}
% ================================================================== %

The master switch \code{TNNPNumericalProtection} (default \code{false})
gates a local guard invoked at every external entry into the embedded TNNP
kernel (\code{calculateCurrent}, \code{solveODE\_start},
\code{solveODE\_end}, \code{derivatives}).

\begin{center}
\begin{tabular}{@{}lll@{}}
\toprule
Key & Type & Default \\
\midrule
\code{TNNPNumericalProtection} & Switch & \code{false} \\
\code{TNNPClampGates} & Switch & \code{true} \\
\code{TNNPClampVmForRates} & Switch & \code{true} \\
\code{TNNPLogNumericalProtection} & Switch & \code{true} \\
\code{TNNPConcentrationFloor} & scalar (mM) & $10^{-12}$ \\
\code{TNNPVMinForRates} & scalar (mV) & $-200$ \\
\code{TNNPVMaxForRates} & scalar (mV) & $+100$ \\
\code{TNNPVoltageZeroEps} & scalar (mV) & $10^{-8}$ \\
\code{TNNPMaxProtectionLogEntries} & label & $10^{6}$ \\
\code{TNNPProtectionLog} & word & \code{TNNP\_numericalProtection\_summary.dat} \\
\bottomrule
\end{tabular}
\end{center}

When protection is enabled, the guards applied in order are:
\begin{enumerate}[topsep=0pt,itemsep=2pt]
  \item \textbf{Non-finite $V$}: \code{!std::isfinite(V)} $\rightarrow$
        reset to $-86.2$ mV.
  \item \textbf{$V$ clamp for rates} (if \code{TNNPClampVmForRates}):
        clamp to $[\,$\code{TNNPVMinForRates}, \code{TNNPVMaxForRates}$\,]$.
  \item \textbf{$V=0$ singularity avoidance}: push $V$ off zero by
        \code{TNNPVoltageZeroEps} (sign preserving) when
        $\lvert V \rvert$ falls below that threshold.
  \item \textbf{Concentration floor} on $K_{i}, Na_{i}, Ca_{i}, Ca_{SR}$:
        non-finite or below-floor entries are reset to
        \code{TNNPConcentrationFloor}.
  \item \textbf{Gate clamp} on $Xr_{1}, Xr_{2}, Xs, m, h, j, d, f, fCa,
        s, r, g$ (if \code{TNNPClampGates}): non-finite to zero,
        clamp to $[0, 1]$.
\end{enumerate}

The embedded implementation stores a capped correction count and emits a
compact summary file at exit; it does not depend on external ionic-model
logging utilities.

% ================================================================== %
\section{Activation tracking and benchmarks}
% ================================================================== %

\paragraph{Activation time per cell.}
First-crossing of $V_{\text{thr}}$ in a time step is recorded by linear
interpolation:
\begin{equation}
  w \;=\; \mathrm{clip}\!\Bigl(
      \frac{V_{\text{thr}} - V^{n}}{V^{n+1} - V^{n}},
      \; 0, 1
    \Bigr),
   \qquad
   t_{\text{act}} \;=\; t^{n} + w\, \Delta t,
\end{equation}
with a per-cell flag cleared after the first crossing so subsequent
re-crossings are ignored.

\paragraph{Early termination on P8.}
When \code{stopAfterPointActivation = true}, the cell nearest
\code{stopActivationPoint} (P8 in the Niederer benchmark) is monitored
and \code{effectiveEndTime} is set to $t_{\text{act}} +
$~\code{stopDelayAfterActivation} once its $t_{\text{act}} > 0$.

\paragraph{Diagonal sampling.}
\code{sampleIDW} performs inverse-distance-weighted interpolation at
\code{nDiagonalSamples} points between \code{diagonalStart} and
\code{diagonalEnd}, with $k = $~\code{sampleNeighbours},
$p = $~\code{samplePower}. Because a diagonal sample averages its $k$ nearest
cells and unactivated cells contribute $t_{\text{act}}=0$, a sample watched
\emph{live} climbs from below as the front fills its neighbourhood; the final
value (all $k$ neighbours activated) is the correct one.

\paragraph{Live activation output.}
\code{writeActivationSamples} is called once before the time loop (so
\code{diagonalActivationTime.dat} and \code{P8\_activationTime.dat} exist from
$t=0$ with every point pending) and again whenever the activated-cell count
increases, overwriting the files. The run can therefore be monitored as the
front advances; the driver exposes the live file in each run directory by a
symbolic link that is replaced by a real copy at the end.

\paragraph{Activation-failure detection.}
When \dictkey{abortOnActivationFailure = true} the solver stops early
(cleanly, exit~0) if the activation fails, avoiding a full run to
\code{endTime}. Two criteria, using the running activated-cell count $N_a$ and
the time $t_\star$ of the last new activation:
\begin{itemize}[topsep=2pt,itemsep=1pt]
  \item \textbf{no ignition:} $N_a=0$ and
        $t>\dictkey{activationIgnitionTimeout}$;
  \item \textbf{conduction block:} $N_a>0$, the stop point is not yet reached,
        and $t-t_\star>\dictkey{activationStallTime}$.
\end{itemize}
The outcome is always recorded in
\code{postProcessing/.../activationStatus.dat} as one of \code{OK},
\code{FAILED\_NO\_IGNITION}, \code{FAILED\_CONDUCTION\_BLOCK} or
\code{INCOMPLETE} (with the final time and $N_a$). The driver reads this file
and, on a failure, logs it to \code{activation\_status\_summary.dat} and
continues with the next configuration rather than aborting the sweep.

% ================================================================== %
\section{Solver flow chart}
% ================================================================== %

The pseudo-code below follows the actual solver.  The shared setup is
common to all nonlinear methods; then one of the three method-specific
blocks is executed according to \dictkey{nonlinearMethod}.

\begin{algorithm}[H]
\caption{Shared time-step setup and final commit/rollback.}
\begin{algorithmic}[1]
\State Read dictionaries and assemble (once) $\Mmass$, $\Kstif$, $\Lop=\Kstif/(\chi C_m)$.
\State Initialise the persistent linear-solver state (PETSc KSP /
       Eigen factorisation) and mark $\Delta t$ as not assembled.
\While{$t^n < T$}
  \State Save snapshots: $\mathsf{V}^n$, $\Iion^n$, $\mathbf{s}^n$.
  \If{$\Delta t$ changed with respect to the assembled one}
    \State (Re)build $\Aimp=\Mmass/\Delta t-\theta\Lop$ and
           $\Bimp=\Mmass/\Delta t+(1-\theta)\Lop$; mark the solver for rebuild.
  \EndIf \Comment{constant if $\Delta t$ is fixed $\Rightarrow$ they are reused (Phase B)}
  \State Build old source $S^n=-\Iion^n+\Iext^n/(\chi C_m)$.
  \If{\dictkey{jfnkInitGuessOrder}$\ge 1$ and previous step exists}
    \State $\mathsf{x}_0\gets\mathsf{V}^n+
      (\Delta t/\Delta t_{\rm prev})(\mathsf{V}^n-\mathsf{V}^{n-1})$.
    \State Clamp $\mathsf{x}_0$ to
      $[\dictkey{jfnkInitGuessVmMin},\dictkey{jfnkInitGuessVmMax}]$.
  \Else
    \State $\mathsf{x}_0\gets\mathsf{V}^n$.
  \EndIf
  \State Run Algorithm~\ref{alg:picard},~\ref{alg:diagonal} or~\ref{alg:jfnk}.
  \If{not converged and \dictkey{nonlinearAcceptUnconverged} is false}
    \State Roll back $\Vm,\Iion,\mathbf{s}$ to the saved $t^n$ snapshots.
  \Else
    \State Commit the final nonlinear iterate.
  \EndIf
  \State Update activation times and benchmark samples.
\EndWhile
\end{algorithmic}
\end{algorithm}

\begin{algorithm}[H]
\caption{Picard nonlinear solve.}
\label{alg:picard}
\begin{algorithmic}[1]
\Require Initial candidate $\mathsf{x}_0$, state snapshot $\mathbf{s}^n$.
\For{$k=0,\dots,\code{nonlinearIterations}-1$}
  \State Store previous iterate $\mathsf{x}_{k-1}$, $I_{\mathrm{ion},k-1}$,
         $\mathbf{s}_{k-1}$ for residuals.
  \State \Call{evaluateNonlinearFields}{$\mathsf{x}_k$}: \Comment{full evaluation \#1}
    reset states to $\mathbf{s}^n$, reconstruct $\Vm$ at Iion quadrature
    points if enabled, call \code{ionicModel->solveODE}, compute $I_{\mathrm{ion},k}$.
  \State Build RHS $\mathsf{r}(I_{\mathrm{ion},k})$ using~\eqref{eq:rhsIion}.
  \State Solve $\Aimp\,\mathsf{V}_*=\mathsf{r}(I_{\mathrm{ion},k})$ with the
    linear back-end, reusing the factorisation if $\Aimp$ did not change
    (\emph{Reuse}; Picard has a constant $\Aimp$).
  \State Relax:
    $\mathsf{x}_{k+1}\gets\mathsf{x}_{k}+
    \code{nonlinearRelaxation}(\mathsf{V}_*-\mathsf{x}_{k})$.
  \State \Call{evaluateNonlinearFields}{$\mathsf{x}_{k+1}$} \Comment{full evaluation \#2 (for residuals)}
  \State Compute $R_{\mathrm{cpl}}$, $R_{\Vm}$, $R_{\Iion}$ and optional
    $R_{\mathbf{s}}$.
  \If{all requested criteria pass and $k+1\ge\code{minNonlinearIterations}$}
    \State Mark converged and exit loop.
  \EndIf
\EndFor
\State \textbf{Note:} each iteration performs \emph{two} full field evaluations
  (LRE reconstruction + ODE + quadrature), which dominate the per-step cost; the
  linear solve is a minor fraction.
\end{algorithmic}
\end{algorithm}

\begin{algorithm}[H]
\caption{\code{diagonalIion} nonlinear solve.}
\label{alg:diagonal}
\begin{algorithmic}[1]
\Require Initial candidate $\mathsf{x}_0$, state snapshot $\mathbf{s}^n$.
\For{$k=0,\dots,\code{nonlinearIterations}-1$}
  \State Evaluate states and $I_{\mathrm{ion},k}$ at $\mathsf{x}_k$ exactly as in Picard.
  \State Approximate the local derivative
    $S'_k\approx-[\Iion(\mathsf{x}_k+\varepsilon_d)-\Iion(\mathsf{x}_k)]/\varepsilon_d$,
    with $\varepsilon_d=\dictkey{diagonalIionEpsilon}$.
  \State Assemble
    $\Aimp_k=\Aimp-\theta\,\mathrm{diag}(S'_k)$ \Comment{only the diagonal changes; the pattern is kept}
  \State Build corrected RHS
    $\mathsf{r}_k=\mathsf{r}(I_{\mathrm{ion},k})-\theta\,\mathrm{diag}(S'_k)\mathsf{x}_k$.
  \State Solve $\Aimp_k\mathsf{V}_*=\mathsf{r}_k$ refreshing only the matrix
    values (\emph{ValuesOnly}); with \code{petscReusePreconditioner} the
    preconditioner is also reused across iterations.
  \State Relax, re-evaluate ODEs/$\Iion$, compute residuals and test convergence
    as in Picard.
\EndFor
\end{algorithmic}
\end{algorithm}

\begin{algorithm}[H]
\caption{Jacobian-Free Newton--Krylov solve.}
\label{alg:jfnk}
\begin{algorithmic}[1]
\Require Initial candidate $\mathsf{x}_0$, state snapshot $\mathbf{s}^n$.
\For{$k=0,\dots,\code{nonlinearIterations}-1$}
  \State Evaluate $\Fres_k=\Fres(\mathsf{x}_k)$; this resets
    states to $\mathbf{s}^n$, solves ODEs and recomputes $\Iion$.
  \State Compute residuals. If all requested criteria pass, mark converged
    and exit.
  \State Define matrix-free product
    $\mathcal{J}_k\mathbf{v}=
    [\Fres(\mathsf{x}_k+\varepsilon\mathbf{v})-\Fres_k]/\varepsilon$,
    with $\varepsilon$ from \dictkey{jfnkEpsilon} \Comment{each matvec = 1 evaluation of $\Fres$ = 1 ODE integration}.
  \If{\dictkey{jfnkPreconditioner} $=$ \code{diffusion}}
    \State Prepare (once / when $\Delta t$ changes) the incomplete ILU of
      $\mathbf{M}=\Aimp$ as right preconditioner $\mathbf{M}^{-1}$.
  \EndIf
  \State Solve $\mathcal{J}_k\delta_k=-\Fres_k$ by
    restarted GMRES (back-end \dictkey{jfnkLinearSolverBackend}: Eigen's in-house GMRES or
    PETSc's \code{KSPGMRES}/\code{MatShell}), using
    \dictkey{jfnkMaxKrylovIterations}, \dictkey{jfnkMaxRestarts},
    \dictkey{jfnkLinearTolerance} and, where relevant, the preconditioner $\mathbf{M}^{-1}$.
  \If{\dictkey{jfnkLineSearch}}
    \State Start $\alpha=\code{nonlinearRelaxation}$ and backtrack until
      Armijo~\eqref{eq:armijo} succeeds, the trial budget is exhausted,
      or \dictkey{jfnkLineSearchAlphaMin} is reached.
  \Else
    \State $\alpha\gets\code{nonlinearRelaxation}$.
  \EndIf
  \State $\mathsf{x}_{k+1}\gets\mathsf{x}_{k}+\alpha\delta_k$.
  \State Re-evaluate $\Fres(\mathsf{x}_{k+1})$, update fields, compute
    residuals and test convergence.
\EndFor
\end{algorithmic}
\end{algorithm}

\noindent\textbf{Legend.} \pyconst{Blue} = Python constant set in
\code{run\_convergence.py}; \dictkey{green} = OpenFOAM dictionary key
read by the solver. Constants and keys are mapped one-to-one in
section~\ref{sec:driver}.

% ================================================================== %
\section{Manufactured driver: \texttt{run\_convergence.py}}
\label{sec:driver}
% ================================================================== %

The convergence script
\code{/home/pablo/OpenFOAM/pablo-v2312/run/tutorials\_electro/highOrderManufacturedFDAImplicitJfNK/run\_convergence.py}
generates an MMS case whose numerical path is intended to match this
solver. The tables below describe the constants exposed there, the
OpenFOAM dictionary entries they generate, and which part of the
discrete equations they affect. The same dictionary keys are read by
\code{highOrderElectroActivationFoamImplicitPETSc}; only problem-specific
items such as exact fields versus TNNP/stimulus differ.

\subsection{Geometry and parameter sweep}
\begin{longtable}{@{}p{0.30\linewidth}p{0.28\linewidth}p{0.36\linewidth}@{}}
\toprule
Python constant & OpenFOAM key & Affected equation / mechanism \\
\midrule \endhead
\code{DIMENSIONS} & --- &
   Selects 2D or 3D manufactured domain and the corresponding LRE
   stencils in \code{P\_LEVELS\_BY\_DIM}. \\
\code{N\_VALUES} & \code{blockMeshDict} cell counts &
   Spatial refinement. This changes FV/LRE discretisation error in the
   PDE, Iion quadrature and state reconstruction. \\
\code{DT\_VALUES} & \code{controlDict.deltaT} &
   Outer PDE time step. Also sets the integration interval for every
   local ODE solve. \\
\code{END\_TIME} & \code{controlDict.endTime} &
   Final time of the transient MMS test. \\
\code{IMPLICIT\_SCHEMES} &
   \code{spatialIntegrationProperties.implicitScheme} &
   Chooses $\theta$ in~\eqref{eq:theta}: \code{backwardEuler}
   ($\theta=1$) or \code{crankNicolson} ($\theta=1/2$). \\
\code{MASS\_MATRIX\_VALUES} &
   \code{spatialIntegrationProperties.massMatrix} &
   Chooses lumped or LRE-consistent $\Mmass$ in
   $\Aimp=\Mmass/\Delta t-\theta\Lop$. \\
\code{ACTIVE\_HO\_PAIRS\_Vm\_States} &
   \code{useHighOrder\_Vm} / \code{useHighOrder\_Iion}
   + \code{LRECoeffs\_*} &
   Tuple \code{(Vm-level, Iion/states-level)}. Vm level affects
   diffusion and consistent mass; Iion/states level affects where ODEs
   are solved and where the nonlinear source is integrated. \\
\code{P\_LEVELS\_BY\_DIM} & \code{highOrderCoeffs} &
   Maps \code{p1,p2,p3} to LRE polynomial order \code{N}, stencil
   target \code{Nn}, radial weight exponent \code{k}, and
   \code{maxStencilSize}. \\
\bottomrule
\end{longtable}

\subsection{Nonlinear method and convergence controls}
\begin{longtable}{@{}p{0.35\linewidth}p{0.30\linewidth}p{0.30\linewidth}@{}}
\toprule
Python constant & OpenFOAM key & Affected equation / mechanism \\
\midrule \endhead
\code{NONLINEAR\_METHOD} & \code{nonlinearMethod} &
   Selects the nonlinear strategy for the PDE/source coupling:
   \code{Picard}, \code{JFNK}, or \code{diagonalIion}. \\
\code{IMPLICIT\_NONLINEAR\_ITERATIONS} & \code{nonlinearIterations} &
   Maximum number of nonlinear corrections per PDE time step. \\
\code{IMPLICIT\_MIN\_NONLINEAR\_ITERATIONS} &
   \code{minNonlinearIterations} &
   Minimum number of corrections before a convergence test may stop the
   loop. \\
\code{NONLINEAR\_TOLERANCE} & \code{nonlinearTolerance} &
   Coupled PDE residual tolerance $R_{\mathrm{cpl}}$. Used directly by
   JFNK and diagonalIion; in the electro solver it is also reported for
   Picard. \\
\code{NONLINEAR\_VM\_TOLERANCE} & \code{nonlinearVmTolerance} &
   Relative correction of the PDE unknown $\Vm$. \\
\code{NONLINEAR\_STATES\_TOLERANCE} &
   \code{nonlinearStatesTolerance} &
   Relative correction of all ODE states. In the MMS solver these are
   $u_1,u_2,u_3$; in this electro solver these are the TNNP state
   variables. \\
\code{NONLINEAR\_REQUIRE\_STATES\_CONVERGENCE} &
   \code{nonlinearRequireStatesConvergence} &
   If true, the state residual is part of the stopping criterion. \\
\code{NONLINEAR\_ACCEPT\_UNCONVERGED} &
   \code{nonlinearAcceptUnconverged} &
   If false, a failed electro nonlinear solve rolls back $\Vm,\Iion$ and
   states to $t^n$. \\
\code{NONLINEAR\_IION\_TOLERANCE} & \code{nonlinearIionTolerance} &
   Relative correction of the nonlinear source $\Iion$. \\
\code{NONLINEAR\_RELAXATION} & \code{nonlinearRelaxation} &
   Picard/diagonal mixing factor or initial Newton step length for
   JFNK/Armijo. \\
\bottomrule
\end{longtable}

\subsection{JFNK and diagonal-Iion controls}
\begin{longtable}{@{}p{0.35\linewidth}p{0.30\linewidth}p{0.30\linewidth}@{}}
\toprule
Python constant & OpenFOAM key & Affected equation / mechanism \\
\midrule \endhead
\code{JFNK\_MAX\_KRYLOV\_ITERATIONS} & \code{jfnkMaxKrylovIterations} &
   Restart length $m$ of GMRES(m), controlling memory and Krylov
   approximation quality. \\
\code{JFNK\_MAX\_RESTARTS} & \code{jfnkMaxRestarts} &
   Number of times the GMRES basis may be discarded and rebuilt. \\
\code{JFNK\_LINEAR\_TOLERANCE} & \code{jfnkLinearTolerance} &
   Relative tolerance of the Newton correction solve
   $\mathbf{J}\delta=-\Fres$. \\
\code{JFNK\_EPSILON} & \code{jfnkEpsilon} &
   Base perturbation for matrix-free Jacobian-vector products. It should
   be above the ODE/noise floor but small enough to approximate a
   derivative. \\
\code{JFNK\_INIT\_GUESS\_ORDER} & \code{jfnkInitGuessOrder} &
   Chooses lagged or linearly extrapolated initial $\Vm$ for the
   nonlinear solve. \\
\code{JFNK\_INIT\_GUESS\_VM\_MIN/MAX} &
   \code{jfnkInitGuessVmMin/Max} &
   Clamp interval for extrapolated guess and, if enabled, ODE inputs. \\
\code{JFNK\_CLAMP\_ODE\_INPUT} & \code{jfnkClampODEInput} &
   Protects the ODE solve from unphysical JFNK finite-difference probes. \\
\code{JFNK\_LINE\_SEARCH} & \code{jfnkLineSearch} &
   Turns Armijo backtracking on/off. \\
\code{JFNK\_LINE\_SEARCH\_MAX\_ITER} &
   \code{jfnkLineSearchMaxIter} &
   Maximum number of halvings of the Newton step. \\
\code{JFNK\_LINE\_SEARCH\_ALPHA\_MIN} &
   \code{jfnkLineSearchAlphaMin} &
   Floor for the Armijo step length. \\
\code{JFNK\_ARMIJO\_C} & \code{jfnkArmijoC} &
   Sufficient-decrease parameter in~\eqref{eq:armijo}. \\
\code{DIAGONAL\_IION\_EPSILON} & \code{diagonalIionEpsilon} &
   Finite-difference step for the diagonal source derivative in
   \code{diagonalIion}. \\
\bottomrule
\end{longtable}

\subsection{Linear PDE solver}
\begin{longtable}{@{}p{0.30\linewidth}p{0.30\linewidth}p{0.35\linewidth}@{}}
\toprule
Python constant & OpenFOAM key & Affected equation / mechanism \\
\midrule \endhead
\code{LINEAR\_SOLVER\_BACKEND} & \code{linearSolverBackend} &
   Chooses the assembled back-end: \code{Eigen} or \code{PETSc}
   (Section~\ref{sec:linalg}). \\
\code{IMPLICIT\_LINEAR\_SOLVER} & \code{linearSolver} &
   Sparse solver used by Picard and diagonalIion for
   $\Aimp\mathsf{V}=\mathsf{r}$ (\code{SparseLU}/\code{BiCGSTAB} in Eigen; mapped to
   KSP/PC in PETSc). JFNK instead uses GMRES on the Jacobian-free shell. \\
\code{PETSC\_LINEAR\_KSP\_TYPE} / \code{PETSC\_LINEAR\_PC\_TYPE} &
   \code{petscLinearKspType} / \code{petscLinearPcType} &
   KSP and preconditioner type when the back-end is PETSc
   (\code{gmres}+\code{ilu}, \code{cg}+\code{gamg}, \dots). \\
\code{IMPLICIT\_TOLERANCE} & \code{tolerance} &
   Relative tolerance of the iterative linear solve. \\
\code{IMPLICIT\_MAX\_ITERATIONS} & \code{maxIterations} &
   Iteration cap of the iterative linear solve. \\
\bottomrule
\end{longtable}

\subsection{State ODE integrator}
\begin{longtable}{@{}p{0.30\linewidth}p{0.30\linewidth}p{0.35\linewidth}@{}}
\toprule
Python constant & OpenFOAM key & Affected equation / mechanism \\
\midrule \endhead
\code{STATE\_ODE\_SOLVER} & \code{stateODESolver} /
   \code{timeIntegrationProperties.solver} &
   Local ODE integrator for MMS states or TNNP states. \\
\code{STATE\_ODE\_ABS\_TOL} & \code{stateODEAbsTol} /
   \code{timeIntegrationProperties.absTol} &
   Absolute tolerance of the ODE solve. \\
\code{STATE\_ODE\_REL\_TOL} & \code{stateODERelTol} /
   \code{timeIntegrationProperties.relTol} &
   Relative ODE tolerance. This controls the noise floor seen by JFNK
   finite differences. \\
\code{STATE\_ODE\_INITIAL\_STEP} & \code{stateODEInitialStep} /
   \code{initialODEStep} &
   First RKF45 sub-step attempted inside one PDE time step. \\
\code{STATE\_ODE\_MAX\_STEPS} & \code{stateODEMaxSteps} /
   \code{maxSteps} &
   Maximum ODE substeps per local solve. \\
\bottomrule
\end{longtable}

\subsection{LRE high-order parameters per p-level}
\code{P\_LEVELS\_BY\_DIM[dim][p]} gives, for each (dim, p):
\code{N} (polynomial order), \code{Nn} (stencil neighbours),
\code{k} (search radius), \code{maxStencilSize} (cap). The script
writes these into a nested
\code{highOrderCoeffs\{LRECoeffs\_Vm\{\dots\}; LRECoeffs\_Iion\{\dots\}\}}
block whenever HO is enabled for either Vm or $\Iion$.

\subsection{Diagnostics and run-control}
\begin{longtable}{@{}p{0.35\linewidth}p{0.30\linewidth}p{0.30\linewidth}@{}}
\toprule
Python constant & OpenFOAM key & Affected mechanism \\
\midrule \endhead
\code{WRITE\_NONLINEAR\_RESIDUALS} &
   \code{writeNonlinearResiduals} & Residual log on/off. \\
\code{KEEP\_LOGS\_ON\_SUCCESS} & --- &
   Python-only switch: keep solver/blockMesh logs after successful runs.
   Failed-run logs are always kept. \\
\bottomrule
\end{longtable}

For the electro-activation tutorial the analogous driver also exposes
activation-threshold, P8 stopping and diagonal-sampling parameters. They
do not affect the PDE/ODE solve itself; they affect only stopping and
post-processing of activation times.

% ================================================================== %
\section{Results}
\label{sec:results}
% ================================================================== %

This section collects the current results for the 2D and 3D
Niederer-style benchmarks. The 2D benchmark (Section~\ref{sec:results:2D}) was re-run
with an extended \code{Picard} sweep covering the state-integration modes and
the front-aware hybrid introduced in
Section~\ref{sec:stateint}; the 3D benchmark (Section~\ref{sec:results:3D}) keeps
the previously reported results (\code{Picard} and \code{diagonalIion}) and is not
modified here. All runs use the Crank--Nicolson time scheme
($\theta = 1/2$), the high-order consistent mass matrix when LRE is
enabled, and the \code{RKF45} state-ODE integrator. Wall-clock time
and peak RSS are defined in Section~\ref{sec:res2d:cost}, together
with the hardware-independent cost quantities.

\subsection{2D Niederer-style benchmark}
\label{sec:results:2D}

The 2D case is the planar $20\times 8$~mm reduction (a single cell in $z$, planar
fibres). The sweep reported here is exhaustive in \code{Picard} and lives under
\path{highOrderNiedererEtAl2012ImplicitPETSc/results/example0/2D_run_20260705_110250/}.
It covers the product of:
\begin{itemize}[topsep=2pt,itemsep=1pt]
  \item \textbf{Mesh:} structured hexahedral (\code{hexa}) and unstructured
        triangular (\code{triangularUnstr}).
  \item \textbf{Hybrid:} \code{frontAwareHybrid} \code{false} (8 configurations) and
        \code{true} (4 configurations).
  \item \textbf{Time step:} $\Delta t\in\{10^{-4},10^{-5},10^{-6}\}$~s.
  \item \textbf{Spatial mesh:} $\Delta x\in\{0.5,0.25,0.125,0.1,0.05\}$~mm. The two
        finest levels ($0.1$ and \SI{0.05}{mm}) were run only for the
        configurations that reach the converged regime; combinations that were not
        run appear as \texttt{--} in the tables.
\end{itemize}
Each configuration is identified by the triplet \code{(Vm, states, Iion)}:
\begin{itemize}[topsep=2pt,itemsep=1pt]
  \item \code{Vm}$\in\{$\code{NO},\code{p3}$\}$: order of the LRE reconstruction of the
        potential in the diffusion operator (\code{NO} = classical finite volumes).
  \item \code{states}$\in\{$\code{na},\code{GP},\code{CCp1},\code{CCp3}$\}$: state
        integration mode (Section~\ref{sec:stateint}). \code{na} = not applicable
        (base case without high order in $\Iion$); \code{GP} = \code{gaussPointODE} (ODE
        at every Gauss point); \code{CCp$k$} = \code{cellCentredReconstruct} (ODE
        at the cell centre + LRE reconstruction of order $k$ of the states to the
        Gauss points).
  \item \code{Iion}$\in\{$\code{NO},\code{p1},\code{p3}$\}$: order of the LRE
        reconstruction at the ionic quadrature points.
\end{itemize}
When \code{frontAwareHybrid} is active, the \code{CCp$k$} configurations
solve the ODE at the Gauss points only in the front cells and by
\code{cellCentredReconstruct} everywhere else (Section~\ref{sec:hybrid}).

\subsubsection{Computational cost quantities}
\label{sec:res2d:cost}

The run time of a simulation depends on which other processes shared the
machine, so it is not a clean measure of the cost of the algorithm (see the
thread-contention analysis in the associated comparison note). To be able to compare
configurations reproducibly, each table reports, besides time and
memory, two \emph{deterministic} quantities (independent of the hardware and of the
machine load):

\begin{description}[topsep=2pt,itemsep=2pt,leftmargin=1.4em]
  \item[wall {[min]}] \emph{Wall-clock time}: elapsed solver time
        (\code{setup}+\code{loop}) reported by its internal timer, rounded to
        minutes. It is what the run ``takes'', but it \textbf{depends on the machine
        load} (CPU contention, thread oversubscription).
  \item[CPU {[min]}] \emph{CPU time} (\code{User time} from
        \texttt{/usr/bin/time -v}): sum of CPU-seconds consumed by the process in
        user mode across all threads. It is less sensitive to load than the
        wall time, but it still inflates when threads \emph{busy-wait} under
        contention; that is why $\text{CPU}>\text{wall}$ in the multi-threaded runs (e.g.\
        at $\Delta x=\SI{0.125}{mm}$).
  \item[RSS {[MiB]}] \emph{Peak resident set size}: maximum resident memory of the
        process, $\lfloor\text{kB}/1024\rfloor$. It scales with the LRE stencil size
        and with the number of cells.
  \item[$N_{\mathrm{nl}}$] \emph{Total nonlinear iterations}: sum over all
        time steps of the number of \code{Picard} iterations. It equals the
        number of full nonlinear residual evaluations --- each one integrates
        the ODEs, evaluates $\Iion$ and assembles/solves the linear PDE. It is
        \textbf{deterministic} (given $\Delta t$, the mesh and the tolerances) and
        therefore the cleanest cost reference across configurations.
  \item[$N_{\mathrm{lin}}$] \emph{Total linear iterations}: sum of the
        GMRES iterations across all linear PDE solves.
        Deterministic; it measures the work of the linear solver.
\end{description}

\noindent\textbf{Important warning about $N_{\mathrm{nl}}$ and $N_{\mathrm{lin}}$.}
These two count iterations of the \emph{solver}, not the cost of integrating the ODE at
each evaluation. That per-evaluation cost differs strongly between modes:
\code{gaussPointODE} integrates $N_q$ stiff ODEs per cell (one per Gauss point),
whereas \code{cellCentredReconstruct} and the smooth cells of the hybrid integrate
\emph{one} per cell. This is why two runs with the same $N_{\mathrm{nl}}$ can
differ several-fold in wall/CPU (see the hybrid discussion below): the
wall and CPU times capture the cost of integrating the ODEs, whereas
$N_{\mathrm{nl}}$ and $N_{\mathrm{lin}}$ capture the number of solves. The four
quantities are complementary.

The number of time steps $n_{\text{steps}}$ (also deterministic) is not tabulated
separately: for $\Delta t\le 10^{-5}$~s \code{Picard} typically converges in one
iteration per step, so that $N_{\mathrm{nl}}\approx n_{\text{steps}}$; at
$\Delta t=10^{-4}$~s several iterations per step are needed and
$N_{\mathrm{nl}}>n_{\text{steps}}$. Runs stop as soon as P8 (the far
corner) activates, so $n_{\text{steps}}\approx t_{\text{P8}}/\Delta t$; the
runs marked \textsc{fail} are aborted at $t=\SI{0.02}{s}$ (ignition timeout).

%%%% HEXA HybridFalse
\subsubsection{Structured hexahedral mesh, no hybrid (\code{frontAwareHybrid false})}
\label{sec:res2d:hexaF}
\begin{center}\footnotesize
\begin{longtable}{l|lrr|rrr|rrr}
\caption{Per-case performance --- structured hexahedral mesh, no hybrid, all \code{Picard}. For each triplet \code{(Vm, states, Iion)} the rows list the available $\Delta x$ [mm] and the columns the three $\Delta t$, grouped three metrics at a time. \textsc{fail} = no ignition (the run is aborted at $t=\SI{0.02}{s}$ by the failure detector); a P8 value implies full activation of the domain; \texttt{--} = combination not run. $N_{\mathrm{nl}}$ and $N_{\mathrm{lin}}$ are the totals of nonlinear and linear iterations (deterministic); wall and CPU in minutes; RSS in MiB.}\label{tab:res2d:hexaF}\\
\toprule $ \Delta t [\text{s}]$ & $10^{-4}$ & $10^{-5}$ & $10^{-6}$ & $10^{-4}$ & $10^{-5}$ & $10^{-6}$ & $10^{-4}$ & $10^{-5}$ & $10^{-6}$\\ \midrule\endfirsthead
\toprule $ \Delta t [\text{s}]$ & $10^{-4}$ & $10^{-5}$ & $10^{-6}$ & $10^{-4}$ & $10^{-5}$ & $10^{-6}$ & $10^{-4}$ & $10^{-5}$ & $10^{-6}$\\ \midrule\endhead
\multicolumn{10}{c}{\textbf{\code{(Vm, states, Iion) = (NO, na, NO)}}}\\\hline
$ \Delta x [\text{mm}] $ & \multicolumn{3}{c|}{P8 [ms]} & \multicolumn{3}{c|}{$N_{\mathrm{nl}}$} & \multicolumn{3}{c|}{$N_{\mathrm{lin}}$}\\
0.5 & 153.32 & 157.15 & 157.31 & 6\,293 & 16\,136 & 161\,519 & 18\,879 & 32\,272 & 323\,038\\
0.25 & 62.99 & 67.30 & 67.42 & 3\,510 & 6\,854 & 68\,652 & 12\,748 & 13\,708 & 137\,304\\
0.125 & 45.03 & 50.05 & 50.26 & 2\,237 & 5\,078 & 50\,971 & 11\,183 & 15\,233 & 101\,942\\
0.1 & 42.76 & 47.86 & 48.11 & 2\,067 & 4\,851 & 48\,751 & 11\,573 & 14\,553 & 97\,502\\
0.05 & 40.22 & 45.28 & 45.60 & 1\,618 & 4\,583 & 46\,136 & 14\,560 & 18\,331 & 138\,407\\
\hline
$ \Delta x [\text{mm}] $ & \multicolumn{3}{c|}{wall [min]} & \multicolumn{3}{c|}{CPU [min]} & \multicolumn{3}{c|}{RSS [MiB]}\\
0.5 & 13 & 27 & 66 & 13 & 27 & 66 & 181 & 181 & 181\\
0.25 & 27 & 46 & 115 & 27 & 46 & 115 & 187 & 188 & 189\\
0.125 & 69 & 131 & 339 & 106 & 208 & 1098 & 213 & 212 & 212\\
0.1 & 104 & 201 & 539 & 150 & 302 & 1546 & 232 & 232 & 232\\
0.05 & 325 & 739 & 1855 & 362 & 842 & 2835 & 393 & 395 & 388\\
\hline
\hline
\multicolumn{10}{c}{\textbf{\code{(Vm, states, Iion) = (NO, GP, p1)}}}\\\hline
$ \Delta x [\text{mm}] $ & \multicolumn{3}{c|}{P8 [ms]} & \multicolumn{3}{c|}{$N_{\mathrm{nl}}$} & \multicolumn{3}{c|}{$N_{\mathrm{lin}}$}\\
0.5 & 91.05 & 94.34 & 94.48 & 4\,312 & 9\,722 & 97\,348 & 12\,936 & 19\,444 & 194\,696\\
0.25 & 54.86 & 58.85 & 58.95 & 2\,829 & 5\,998 & 60\,074 & 10\,126 & 11\,996 & 120\,148\\
0.125 & 43.21 & 48.04 & 48.27 & 1\,886 & 4\,875 & 48\,965 & 9\,430 & 14\,625 & 97\,930\\
0.1 & 41.84 & 46.79 & 47.06 & 1\,690 & 4\,743 & 47\,688 & 9\,406 & 14\,229 & 95\,376\\
0.05 & 40.14 & 45.18 & 45.50 & 1\,611 & 4\,573 & 46\,036 & 14\,497 & 18\,290 & 138\,108\\
\hline
$ \Delta x [\text{mm}] $ & \multicolumn{3}{c|}{wall [min]} & \multicolumn{3}{c|}{CPU [min]} & \multicolumn{3}{c|}{RSS [MiB]}\\
0.5 & 18 & 35 & 78 & 18 & 35 & 78 & 181 & 182 & 181\\
0.25 & 47 & 85 & 192 & 47 & 85 & 192 & 191 & 193 & 192\\
0.125 & 114 & 254 & 603 & 144 & 334 & 1322 & 228 & 228 & 229\\
0.1 & 174 & 406 & 996 & 211 & 503 & 1975 & 258 & 258 & 257\\
0.05 & 656 & 1535 & 3652 & 694 & 1637 & 4613 & 490 & 492 & 490\\
\hline
\hline
\multicolumn{10}{c}{\textbf{\code{(Vm, states, Iion) = (NO, CCp1, p1)}}}\\\hline
$ \Delta x [\text{mm}] $ & \multicolumn{3}{c|}{P8 [ms]} & \multicolumn{3}{c|}{$N_{\mathrm{nl}}$} & \multicolumn{3}{c|}{$N_{\mathrm{lin}}$}\\
0.5 & 80.23 & \textsc{fail} & \textsc{fail} & 4\,079 & 2\,004 & 20\,002 & 12\,237 & 4\,008 & 40\,004\\
0.25 & 60.37 & \textsc{fail} & \textsc{fail} & 3\,428 & 2\,004 & 20\,002 & 11\,872 & 4\,008 & 40\,004\\
0.125 & 49.71 & \textsc{fail} & \textsc{fail} & 2\,470 & 2\,003 & 20\,002 & 12\,350 & 6\,009 & 40\,004\\
0.1 & 48.06 & \textsc{fail} & \textsc{fail} & 2\,284 & 2\,003 & 20\,002 & 12\,225 & 6\,009 & 40\,004\\
0.05 & 45.95 & \textsc{fail} & \textsc{fail} & 1\,845 & 2\,003 & 20\,002 & 16\,603 & 8\,012 & 60\,006\\
\hline
$ \Delta x [\text{mm}] $ & \multicolumn{3}{c|}{wall [min]} & \multicolumn{3}{c|}{CPU [min]} & \multicolumn{3}{c|}{RSS [MiB]}\\
0.5 & 11 & 6 & 12 & 11 & 6 & 12 & 183 & 183 & 183\\
0.25 & 36 & 23 & 48 & 36 & 23 & 48 & 194 & 196 & 196\\
0.125 & 94 & 80 & 185 & 136 & 111 & 485 & 241 & 241 & 241\\
0.1 & 148 & 135 & 320 & 199 & 177 & 727 & 277 & 276 & 275\\
0.05 & 473 & 530 & 1230 & 516 & 575 & 1644 & 567 & 565 & 564\\
\hline
\hline
\multicolumn{10}{c}{\textbf{\code{(Vm, states, Iion) = (p3, GP, NO)}}}\\\hline
$ \Delta x [\text{mm}] $ & \multicolumn{3}{c|}{P8 [ms]} & \multicolumn{3}{c|}{$N_{\mathrm{nl}}$} & \multicolumn{3}{c|}{$N_{\mathrm{lin}}$}\\
0.5 & 180.55 & 184.69 & 184.90 & 7\,027 & 18\,725 & 187\,450 & 21\,081 & 56\,175 & 562\,350\\
0.25 & 68.84 & 73.57 & 73.73 & 4\,072 & 7\,453 & 74\,677 & 12\,216 & 22\,359 & 224\,031\\
0.125 & 48.83 & 54.36 & 54.58 & 2\,688 & 5\,503 & 55\,232 & 8\,064 & 16\,437 & 165\,696\\
\hline
$ \Delta x [\text{mm}] $ & \multicolumn{3}{c|}{wall [min]} & \multicolumn{3}{c|}{CPU [min]} & \multicolumn{3}{c|}{RSS [MiB]}\\
0.5 & 15 & 38 & 169 & 60 & 211 & 2166 & 192 & 192 & 188\\
0.25 & 31 & 50 & 127 & 58 & 109 & 596 & 227 & 228 & 223\\
0.125 & 81 & 144 & 398 & 126 & 237 & 1357 & 373 & 373 & 373\\
\hline
\hline
\multicolumn{10}{c}{\textbf{\code{(Vm, states, Iion) = (p3, GP, p1)}}}\\\hline
$ \Delta x [\text{mm}] $ & \multicolumn{3}{c|}{P8 [ms]} & \multicolumn{3}{c|}{$N_{\mathrm{nl}}$} & \multicolumn{3}{c|}{$N_{\mathrm{lin}}$}\\
0.5 & 68.92 & 71.41 & 71.53 & 4\,056 & 9\,550 & 95\,641 & 12\,168 & 28\,650 & 286\,923\\
0.25 & 58.81 & 63.12 & 63.27 & 3\,120 & 6\,391 & 64\,047 & 9\,360 & 19\,173 & 192\,141\\
0.125 & 46.33 & 51.39 & 51.61 & 2\,216 & 5\,204 & 52\,246 & 6\,648 & 15\,599 & 156\,738\\
\hline
$ \Delta x [\text{mm}] $ & \multicolumn{3}{c|}{wall [min]} & \multicolumn{3}{c|}{CPU [min]} & \multicolumn{3}{c|}{RSS [MiB]}\\
0.5 & 20 & 38 & 138 & 69 & 155 & 1295 & 193 & 193 & 193\\
0.25 & 52 & 96 & 260 & 91 & 174 & 1038 & 231 & 231 & 231\\
0.125 & 132 & 273 & 689 & 170 & 366 & 1581 & 385 & 385 & 383\\
\hline
\hline
\multicolumn{10}{c}{\textbf{\code{(Vm, states, Iion) = (p3, GP, p3)}}}\\\hline
$ \Delta x [\text{mm}] $ & \multicolumn{3}{c|}{P8 [ms]} & \multicolumn{3}{c|}{$N_{\mathrm{nl}}$} & \multicolumn{3}{c|}{$N_{\mathrm{lin}}$}\\
0.5 & 33.09 & 36.31 & 36.56 & 1\,129 & 3\,717 & 37\,404 & 3\,387 & 11\,151 & 112\,212\\
0.25 & 40.22 & 44.20 & 44.43 & 1\,369 & 4\,494 & 45\,160 & 4\,107 & 13\,482 & 135\,480\\
0.125 & 40.84 & 45.40 & 45.66 & 1\,392 & 4\,604 & 46\,293 & 4\,176 & 13\,812 & 138\,879\\
\hline
$ \Delta x [\text{mm}] $ & \multicolumn{3}{c|}{wall [min]} & \multicolumn{3}{c|}{CPU [min]} & \multicolumn{3}{c|}{RSS [MiB]}\\
0.5 & 21 & 56 & 141 & 35 & 101 & 595 & 198 & 199 & 199\\
0.25 & 96 & 255 & 622 & 113 & 310 & 1171 & 252 & 252 & 252\\
0.125 & 335 & 932 & 2228 & 359 & 1012 & 3026 & 467 & 461 & 463\\
\hline
\hline
\multicolumn{10}{c}{\textbf{\code{(Vm, states, Iion) = (p3, CCp1, p1)}}}\\\hline
$ \Delta x [\text{mm}] $ & \multicolumn{3}{c|}{P8 [ms]} & \multicolumn{3}{c|}{$N_{\mathrm{nl}}$} & \multicolumn{3}{c|}{$N_{\mathrm{lin}}$}\\
0.5 & 91.40 & \textsc{fail} & \textsc{fail} & 4\,636 & 2\,004 & 20\,002 & 13\,908 & 6\,012 & 60\,006\\
0.25 & 66.46 & \textsc{fail} & \textsc{fail} & 4\,168 & 2\,004 & 20\,002 & 12\,504 & 6\,012 & 60\,006\\
0.125 & 52.98 & \textsc{fail} & \textsc{fail} & 2\,925 & 2\,003 & 20\,002 & 8\,775 & 6\,009 & 60\,006\\
0.1 & 50.47 & \textsc{fail} & \textsc{fail} & 2\,554 & 2\,003 & 20\,002 & 10\,214 & 6\,009 & 60\,006\\
0.05 & 46.51 & \textsc{fail} & \textsc{fail} & 1\,879 & 2\,003 & 20\,002 & 9\,393 & 6\,009 & 60\,006\\
\hline
$ \Delta x [\text{mm}] $ & \multicolumn{3}{c|}{wall [min]} & \multicolumn{3}{c|}{CPU [min]} & \multicolumn{3}{c|}{RSS [MiB]}\\
0.5 & 15 & 7 & 24 & 71 & 31 & 267 & 196 & 195 & 194\\
0.25 & 46 & 24 & 62 & 96 & 48 & 304 & 237 & 236 & 238\\
0.125 & 111 & 82 & 201 & 161 & 116 & 544 & 413 & 413 & 413\\
0.1 & 160 & 133 & 322 & 217 & 175 & 739 & 538 & 534 & 533\\
0.05 & 470 & 523 & 1270 & 513 & 569 & 1717 & 1433 & 1416 & 1416\\
\hline
\hline
\multicolumn{10}{c}{\textbf{\code{(Vm, states, Iion) = (p3, CCp3, p3)}}}\\\hline
$ \Delta x [\text{mm}] $ & \multicolumn{3}{c|}{P8 [ms]} & \multicolumn{3}{c|}{$N_{\mathrm{nl}}$} & \multicolumn{3}{c|}{$N_{\mathrm{lin}}$}\\
0.5 & 66.63 & \textsc{fail} & \textsc{fail} & 2\,788 & 2\,004 & 20\,002 & 8\,364 & 6\,012 & 60\,006\\
0.25 & 60.80 & \textsc{fail} & \textsc{fail} & 2\,942 & 2\,004 & 20\,002 & 8\,826 & 6\,012 & 60\,006\\
0.125 & 52.20 & \textsc{fail} & \textsc{fail} & 2\,455 & 2\,003 & 20\,002 & 7\,365 & 6\,009 & 60\,006\\
0.1 & 50.15 & \textsc{fail} & -- & 2\,231 & 2\,003 & -- & 8\,922 & 6\,009 & --\\
0.05 & 46.48 & -- & -- & 1\,863 & -- & -- & 9\,313 & -- & --\\
\hline
$ \Delta x [\text{mm}] $ & \multicolumn{3}{c|}{wall [min]} & \multicolumn{3}{c|}{CPU [min]} & \multicolumn{3}{c|}{RSS [MiB]}\\
0.5 & 10 & 8 & 33 & 44 & 32 & 273 & 203 & 203 & 203\\
0.25 & 37 & 28 & 101 & 73 & 52 & 344 & 271 & 270 & 271\\
0.125 & 107 & 97 & 374 & 149 & 131 & 752 & 541 & 538 & 538\\
0.1 & 163 & 161 & -- & 211 & 202 & -- & 743 & 741 & --\\
0.05 & 540 & -- & -- & 581 & -- & -- & 2416 & -- & --\\
\hline
\bottomrule\end{longtable}\end{center}
%%%% HEXA HybridTrue
\subsubsection{Structured hexahedral mesh, front-aware hybrid (\code{frontAwareHybrid true})}
\label{sec:res2d:hexaT}
\begin{center}\footnotesize
\begin{longtable}{l|lrr|rrr|rrr}
\caption{Per-case performance --- structured hexahedral mesh, front-aware hybrid, all \code{Picard}. For each triplet \code{(Vm, states, Iion)} the rows list the available $\Delta x$ [mm] and the columns the three $\Delta t$, grouped three metrics at a time. \textsc{fail} = no ignition (the run is aborted at $t=\SI{0.02}{s}$ by the failure detector); a P8 value implies full activation of the domain; \texttt{--} = combination not run. $N_{\mathrm{nl}}$ and $N_{\mathrm{lin}}$ are the totals of nonlinear and linear iterations (deterministic); wall and CPU in minutes; RSS in MiB.}\label{tab:res2d:hexaT}\\
\toprule $ \Delta t [\text{s}]$ & $10^{-4}$ & $10^{-5}$ & $10^{-6}$ & $10^{-4}$ & $10^{-5}$ & $10^{-6}$ & $10^{-4}$ & $10^{-5}$ & $10^{-6}$\\ \midrule\endfirsthead
\toprule $ \Delta t [\text{s}]$ & $10^{-4}$ & $10^{-5}$ & $10^{-6}$ & $10^{-4}$ & $10^{-5}$ & $10^{-6}$ & $10^{-4}$ & $10^{-5}$ & $10^{-6}$\\ \midrule\endhead
\multicolumn{10}{c}{\textbf{\code{(Vm, states, Iion) = (NO, na, NO)}}}\\\hline
$ \Delta x [\text{mm}] $ & \multicolumn{3}{c|}{P8 [ms]} & \multicolumn{3}{c|}{$N_{\mathrm{nl}}$} & \multicolumn{3}{c|}{$N_{\mathrm{lin}}$}\\
0.5 & 153.32 & 157.15 & 157.31 & 6\,293 & 16\,136 & 161\,519 & 18\,879 & 32\,272 & 323\,038\\
0.25 & 62.99 & 67.30 & 67.42 & 3\,510 & 6\,854 & 68\,652 & 12\,748 & 13\,708 & 137\,304\\
0.125 & 45.03 & 50.05 & 50.26 & 2\,237 & 5\,078 & 50\,971 & 11\,183 & 15\,233 & 101\,942\\
\hline
$ \Delta x [\text{mm}] $ & \multicolumn{3}{c|}{wall [min]} & \multicolumn{3}{c|}{CPU [min]} & \multicolumn{3}{c|}{RSS [MiB]}\\
0.5 & 13 & 27 & 69 & 13 & 27 & 69 & 181 & 181 & 181\\
0.25 & 28 & 47 & 117 & 28 & 47 & 117 & 188 & 189 & 188\\
0.125 & 69 & 131 & 341 & 108 & 211 & 1119 & 212 & 212 & 212\\
\hline
\hline
\multicolumn{10}{c}{\textbf{\code{(Vm, states, Iion) = (NO, CCp1, p1)}}}\\\hline
$ \Delta x [\text{mm}] $ & \multicolumn{3}{c|}{P8 [ms]} & \multicolumn{3}{c|}{$N_{\mathrm{nl}}$} & \multicolumn{3}{c|}{$N_{\mathrm{lin}}$}\\
0.5 & 91.05 & 94.34 & 94.48 & 4\,320 & 9\,722 & 97\,350 & 12\,960 & 19\,444 & 194\,699\\
0.25 & 54.86 & 58.85 & 58.95 & 2\,826 & 5\,998 & 60\,076 & 10\,117 & 11\,996 & 120\,152\\
0.125 & 43.20 & 48.05 & 48.29 & 1\,904 & 4\,876 & 48\,988 & 9\,519 & 14\,627 & 97\,976\\
0.1 & 41.85 & 46.79 & 47.07 & 1\,695 & 4\,744 & 47\,702 & 9\,430 & 14\,232 & 95\,403\\
0.05 & 40.14 & 45.18 & 45.51 & 1\,611 & 4\,573 & 46\,040 & 14\,497 & 18\,291 & 138\,119\\
\hline
$ \Delta x [\text{mm}] $ & \multicolumn{3}{c|}{wall [min]} & \multicolumn{3}{c|}{CPU [min]} & \multicolumn{3}{c|}{RSS [MiB]}\\
0.5 & 12 & 23 & 57 & 12 & 23 & 57 & 182 & 183 & 183\\
0.25 & 27 & 50 & 126 & 27 & 50 & 126 & 195 & 196 & 195\\
0.125 & 63 & 136 & 385 & 96 & 213 & 1136 & 240 & 241 & 240\\
0.1 & 94 & 220 & 630 & 132 & 319 & 1609 & 278 & 279 & 277\\
0.05 & 340 & 825 & 2208 & 379 & 927 & 3175 & 564 & 568 & 570\\
\hline
\hline
\multicolumn{10}{c}{\textbf{\code{(Vm, states, Iion) = (p3, CCp1, p1)}}}\\\hline
$ \Delta x [\text{mm}] $ & \multicolumn{3}{c|}{P8 [ms]} & \multicolumn{3}{c|}{$N_{\mathrm{nl}}$} & \multicolumn{3}{c|}{$N_{\mathrm{lin}}$}\\
0.5 & 68.93 & 71.44 & 71.56 & 4\,041 & 9\,554 & 95\,681 & 12\,123 & 28\,662 & 287\,043\\
0.25 & 58.82 & 63.11 & 63.27 & 3\,120 & 6\,390 & 64\,044 & 9\,360 & 19\,170 & 192\,132\\
0.125 & 46.33 & 51.39 & 51.61 & 2\,218 & 5\,204 & 52\,250 & 6\,654 & 15\,598 & 156\,750\\
\hline
$ \Delta x [\text{mm}] $ & \multicolumn{3}{c|}{wall [min]} & \multicolumn{3}{c|}{CPU [min]} & \multicolumn{3}{c|}{RSS [MiB]}\\
0.5 & 13 & 27 & 114 & 62 & 143 & 1277 & 195 & 195 & 195\\
0.25 & 31 & 56 & 175 & 69 & 134 & 955 & 238 & 238 & 237\\
0.125 & 72 & 149 & 452 & 111 & 240 & 1371 & 414 & 414 & 414\\
\hline
\hline
\multicolumn{10}{c}{\textbf{\code{(Vm, states, Iion) = (p3, CCp3, p3)}}}\\\hline
$ \Delta x [\text{mm}] $ & \multicolumn{3}{c|}{P8 [ms]} & \multicolumn{3}{c|}{$N_{\mathrm{nl}}$} & \multicolumn{3}{c|}{$N_{\mathrm{lin}}$}\\
0.5 & 33.09 & 36.30 & 36.51 & 1\,130 & 3\,715 & 37\,350 & 3\,390 & 11\,145 & 112\,050\\
0.25 & 40.22 & 44.22 & 44.54 & 1\,371 & 4\,496 & 45\,266 & 4\,113 & 13\,488 & 135\,798\\
0.125 & 40.84 & 45.43 & 45.90 & 1\,391 & 4\,607 & 46\,527 & 4\,173 & 13\,821 & 139\,581\\
0.1 & 40.72 & 45.48 & 45.98 & 1\,402 & 4\,609 & 46\,573 & 5\,607 & 13\,827 & 139\,719\\
0.05 & 40.05 & 45.04 & 45.52 & 1\,604 & 4\,558 & 46\,051 & 8\,019 & 13\,674 & 138\,153\\
\hline
$ \Delta x [\text{mm}] $ & \multicolumn{3}{c|}{wall [min]} & \multicolumn{3}{c|}{CPU [min]} & \multicolumn{3}{c|}{RSS [MiB]}\\
0.5 & 6 & 17 & 64 & 20 & 62 & 519 & 202 & 202 & 202\\
0.25 & 20 & 60 & 216 & 37 & 115 & 767 & 267 & 267 & 267\\
0.125 & 64 & 205 & 845 & 89 & 298 & 1841 & 530 & 530 & 530\\
0.1 & 104 & 306 & 1178 & 135 & 403 & 2147 & 726 & 725 & 726\\
0.05 & 419 & 1123 & 4211 & 456 & 1224 & 5239 & 2370 & 2367 & 2351\\
\hline
\bottomrule\end{longtable}\end{center}
%%%% TRI HybridFalse
\subsubsection{Unstructured triangular mesh, no hybrid}
\label{sec:res2d:triF}
\begin{center}\footnotesize
\begin{longtable}{l|lrr|rrr|rrr}
\caption{Per-case performance --- unstructured triangular mesh, no hybrid, all \code{Picard}. For each triplet \code{(Vm, states, Iion)} the rows list the available $\Delta x$ [mm] and the columns the three $\Delta t$, grouped three metrics at a time. \textsc{fail} = no ignition (the run is aborted at $t=\SI{0.02}{s}$ by the failure detector); a P8 value implies full activation of the domain; \texttt{--} = combination not run. $N_{\mathrm{nl}}$ and $N_{\mathrm{lin}}$ are the totals of nonlinear and linear iterations (deterministic); wall and CPU in minutes; RSS in MiB.}\label{tab:res2d:triF}\\
\toprule $ \Delta t [\text{s}]$ & $10^{-4}$ & $10^{-5}$ & $10^{-6}$ & $10^{-4}$ & $10^{-5}$ & $10^{-6}$ & $10^{-4}$ & $10^{-5}$ & $10^{-6}$\\ \midrule\endfirsthead
\toprule $ \Delta t [\text{s}]$ & $10^{-4}$ & $10^{-5}$ & $10^{-6}$ & $10^{-4}$ & $10^{-5}$ & $10^{-6}$ & $10^{-4}$ & $10^{-5}$ & $10^{-6}$\\ \midrule\endhead
\multicolumn{10}{c}{\textbf{\code{(Vm, states, Iion) = (NO, na, NO)}}}\\\hline
$ \Delta x [\text{mm}] $ & \multicolumn{3}{c|}{P8 [ms]} & \multicolumn{3}{c|}{$N_{\mathrm{nl}}$} & \multicolumn{3}{c|}{$N_{\mathrm{lin}}$}\\
0.5 & 55.81 & 59.70 & 59.88 & 3\,000 & 6\,060 & 60\,760 & 12\,000 & 18\,179 & 121\,520\\
0.25 & 45.18 & 50.28 & 50.51 & 2\,231 & 5\,095 & 51\,168 & 15\,615 & 16\,180 & 102\,336\\
0.125 & 40.93 & 46.00 & 46.30 & 1\,662 & 4\,657 & 46\,858 & 18\,281 & 23\,284 & 140\,574\\
\hline
$ \Delta x [\text{mm}] $ & \multicolumn{3}{c|}{wall [min]} & \multicolumn{3}{c|}{CPU [min]} & \multicolumn{3}{c|}{RSS [MiB]}\\
0.5 & 15 & 25 & 61 & 15 & 25 & 61 & 183 & 183 & 183\\
0.25 & 45 & 88 & 222 & 72 & 150 & 848 & 197 & 197 & 198\\
0.125 & 128 & 304 & 808 & 166 & 407 & 1803 & 246 & 247 & 247\\
\hline
\hline
\multicolumn{10}{c}{\textbf{\code{(Vm, states, Iion) = (NO, GP, p1)}}}\\\hline
$ \Delta x [\text{mm}] $ & \multicolumn{3}{c|}{P8 [ms]} & \multicolumn{3}{c|}{$N_{\mathrm{nl}}$} & \multicolumn{3}{c|}{$N_{\mathrm{lin}}$}\\
0.5 & 55.81 & 59.70 & 59.88 & 3\,000 & 6\,060 & 60\,760 & 12\,000 & 18\,179 & 121\,520\\
0.25 & 45.18 & 50.28 & 50.51 & 2\,231 & 5\,095 & 51\,168 & 15\,615 & 16\,180 & 102\,336\\
0.125 & 40.93 & 46.00 & 46.30 & 1\,662 & 4\,657 & 46\,858 & 18\,281 & 23\,284 & 140\,574\\
\hline
$ \Delta x [\text{mm}] $ & \multicolumn{3}{c|}{wall [min]} & \multicolumn{3}{c|}{CPU [min]} & \multicolumn{3}{c|}{RSS [MiB]}\\
0.5 & 15 & 25 & 60 & 15 & 25 & 60 & 183 & 184 & 183\\
0.25 & 45 & 88 & 231 & 73 & 150 & 853 & 197 & 197 & 198\\
0.125 & 118 & 297 & 687 & 147 & 389 & 1482 & 249 & 248 & 250\\
\hline
\hline
\multicolumn{10}{c}{\textbf{\code{(Vm, states, Iion) = (NO, CCp1, p1)}}}\\\hline
$ \Delta x [\text{mm}] $ & \multicolumn{3}{c|}{P8 [ms]} & \multicolumn{3}{c|}{$N_{\mathrm{nl}}$} & \multicolumn{3}{c|}{$N_{\mathrm{lin}}$}\\
0.5 & 55.81 & 59.70 & 59.88 & 3\,000 & 6\,060 & 60\,760 & 12\,000 & 18\,179 & 121\,520\\
0.25 & 45.18 & 50.28 & 50.51 & 2\,231 & 5\,095 & 51\,168 & 15\,615 & 16\,180 & 102\,336\\
0.125 & 40.93 & 46.00 & 46.30 & 1\,662 & 4\,657 & 46\,858 & 18\,281 & 23\,284 & 140\,574\\
\hline
$ \Delta x [\text{mm}] $ & \multicolumn{3}{c|}{wall [min]} & \multicolumn{3}{c|}{CPU [min]} & \multicolumn{3}{c|}{RSS [MiB]}\\
0.5 & 14 & 25 & 64 & 14 & 25 & 64 & 185 & 185 & 185\\
0.25 & 43 & 86 & 243 & 71 & 149 & 867 & 203 & 203 & 204\\
0.125 & 131 & 320 & 874 & 169 & 421 & 1891 & 272 & 271 & 273\\
\hline
\hline
\multicolumn{10}{c}{\textbf{\code{(Vm, states, Iion) = (p3, GP, NO)}}}\\\hline
$ \Delta x [\text{mm}] $ & \multicolumn{3}{c|}{P8 [ms]} & \multicolumn{3}{c|}{$N_{\mathrm{nl}}$} & \multicolumn{3}{c|}{$N_{\mathrm{lin}}$}\\
0.5 & 75.42 & 57.19 & 57.29 & 4\,320 & 7\,956 & 79\,678 & 12\,960 & 23\,868 & 239\,034\\
0.25 & 51.56 & 56.76 & 56.99 & 2\,981 & 5\,741 & 57\,626 & 10\,672 & 17\,223 & 172\,878\\
0.125 & 43.26 & 48.42 & 48.66 & 2\,144 & 4\,904 & 49\,263 & 10\,719 & 14\,712 & 147\,789\\
\hline
$ \Delta x [\text{mm}] $ & \multicolumn{3}{c|}{wall [min]} & \multicolumn{3}{c|}{CPU [min]} & \multicolumn{3}{c|}{RSS [MiB]}\\
0.5 & 20 & 31 & 88 & 50 & 81 & 608 & 203 & 202 & 200\\
0.25 & 51 & 85 & 217 & 71 & 122 & 584 & 272 & 272 & 268\\
0.125 & 150 & 304 & 770 & 189 & 398 & 1645 & 561 & 560 & 562\\
\hline
\hline
\multicolumn{10}{c}{\textbf{\code{(Vm, states, Iion) = (p3, GP, p1)}}}\\\hline
$ \Delta x [\text{mm}] $ & \multicolumn{3}{c|}{P8 [ms]} & \multicolumn{3}{c|}{$N_{\mathrm{nl}}$} & \multicolumn{3}{c|}{$N_{\mathrm{lin}}$}\\
0.5 & 75.42 & 57.19 & 57.29 & 4\,320 & 7\,956 & 79\,678 & 12\,960 & 23\,868 & 239\,034\\
0.25 & 51.56 & 56.76 & 56.99 & 2\,981 & 5\,741 & 57\,626 & 10\,672 & 17\,223 & 172\,878\\
0.125 & 43.26 & 48.42 & 48.66 & 2\,144 & 4\,904 & 49\,263 & 10\,719 & 14\,712 & 147\,789\\
\hline
$ \Delta x [\text{mm}] $ & \multicolumn{3}{c|}{wall [min]} & \multicolumn{3}{c|}{CPU [min]} & \multicolumn{3}{c|}{RSS [MiB]}\\
0.5 & 23 & 37 & 124 & 75 & 134 & 1090 & 204 & 204 & 204\\
0.25 & 57 & 96 & 269 & 93 & 166 & 971 & 273 & 273 & 273\\
0.125 & 151 & 299 & 804 & 189 & 385 & 1678 & 563 & 561 & 563\\
\hline
\hline
\multicolumn{10}{c}{\textbf{\code{(Vm, states, Iion) = (p3, GP, p3)}}}\\\hline
$ \Delta x [\text{mm}] $ & \multicolumn{3}{c|}{P8 [ms]} & \multicolumn{3}{c|}{$N_{\mathrm{nl}}$} & \multicolumn{3}{c|}{$N_{\mathrm{lin}}$}\\
0.5 & 38.88 & 42.26 & 42.47 & 1\,930 & 4\,291 & 43\,113 & 5\,790 & 12\,873 & 129\,339\\
0.25 & 38.93 & 43.26 & 43.52 & 1\,546 & 4\,388 & 44\,119 & 5\,332 & 13\,164 & 132\,357\\
0.125 & 39.80 & 44.66 & 44.96 & 1\,569 & 4\,526 & 45\,549 & 7\,844 & 13\,578 & 136\,647\\
\hline
$ \Delta x [\text{mm}] $ & \multicolumn{3}{c|}{wall [min]} & \multicolumn{3}{c|}{CPU [min]} & \multicolumn{3}{c|}{RSS [MiB]}\\
0.5 & 40 & 75 & 197 & 64 & 128 & 723 & 211 & 211 & 211\\
0.25 & 127 & 296 & 702 & 146 & 350 & 1240 & 300 & 297 & 300\\
0.125 & 456 & 1069 & 2618 & 486 & 1149 & 3438 & 656 & 665 & 657\\
\hline
\hline
\multicolumn{10}{c}{\textbf{\code{(Vm, states, Iion) = (p3, CCp1, p1)}}}\\\hline
$ \Delta x [\text{mm}] $ & \multicolumn{3}{c|}{P8 [ms]} & \multicolumn{3}{c|}{$N_{\mathrm{nl}}$} & \multicolumn{3}{c|}{$N_{\mathrm{lin}}$}\\
0.5 & 75.42 & 57.19 & 57.29 & 4\,320 & 7\,956 & 79\,678 & 12\,960 & 23\,868 & 239\,034\\
0.25 & 51.56 & 56.76 & 56.99 & 2\,981 & 5\,741 & 57\,626 & 10\,672 & 17\,223 & 172\,878\\
0.125 & 43.26 & 48.42 & 48.66 & 2\,144 & 4\,904 & 49\,263 & 10\,719 & 14\,712 & 147\,789\\
\hline
$ \Delta x [\text{mm}] $ & \multicolumn{3}{c|}{wall [min]} & \multicolumn{3}{c|}{CPU [min]} & \multicolumn{3}{c|}{RSS [MiB]}\\
0.5 & 24 & 38 & 135 & 77 & 136 & 1103 & 205 & 205 & 205\\
0.25 & 59 & 101 & 298 & 95 & 171 & 1000 & 290 & 290 & 291\\
0.125 & 173 & 367 & 988 & 221 & 472 & 2089 & 616 & 619 & 619\\
\hline
\hline
\multicolumn{10}{c}{\textbf{\code{(Vm, states, Iion) = (p3, CCp3, p3)}}}\\\hline
$ \Delta x [\text{mm}] $ & \multicolumn{3}{c|}{P8 [ms]} & \multicolumn{3}{c|}{$N_{\mathrm{nl}}$} & \multicolumn{3}{c|}{$N_{\mathrm{lin}}$}\\
0.5 & 61.74 & \textsc{fail} & \textsc{fail} & 2\,814 & 2\,004 & 20\,002 & 8\,442 & 6\,012 & 60\,006\\
0.25 & 54.04 & \textsc{fail} & \textsc{fail} & 2\,486 & 2\,003 & 20\,002 & 7\,708 & 6\,009 & 60\,006\\
0.125 & 48.77 & \textsc{fail} & \textsc{fail} & 1\,972 & 2\,003 & 20\,002 & 9\,858 & 6\,009 & 60\,006\\
\hline
$ \Delta x [\text{mm}] $ & \multicolumn{3}{c|}{wall [min]} & \multicolumn{3}{c|}{CPU [min]} & \multicolumn{3}{c|}{RSS [MiB]}\\
0.5 & 20 & 15 & 52 & 54 & 40 & 296 & 220 & 219 & 219\\
0.25 & 67 & 59 & 188 & 97 & 84 & 433 & 334 & 334 & 334\\
0.125 & 208 & 230 & 681 & 253 & 275 & 1050 & 791 & 791 & 791\\
\hline
\bottomrule\end{longtable}\end{center}
%%%% TRI HybridTrue
\subsubsection{Unstructured triangular mesh, front-aware hybrid}
\label{sec:res2d:triT}
\begin{center}\footnotesize
\begin{longtable}{l|lrr|rrr|rrr}
\caption{Per-case performance --- unstructured triangular mesh, front-aware hybrid, all \code{Picard}. For each triplet \code{(Vm, states, Iion)} the rows list the available $\Delta x$ [mm] and the columns the three $\Delta t$, grouped three metrics at a time. \textsc{fail} = no ignition (the run is aborted at $t=\SI{0.02}{s}$ by the failure detector); a P8 value implies full activation of the domain; \texttt{--} = combination not run. $N_{\mathrm{nl}}$ and $N_{\mathrm{lin}}$ are the totals of nonlinear and linear iterations (deterministic); wall and CPU in minutes; RSS in MiB.}\label{tab:res2d:triT}\\
\toprule $ \Delta t [\text{s}]$ & $10^{-4}$ & $10^{-5}$ & $10^{-6}$ & $10^{-4}$ & $10^{-5}$ & $10^{-6}$ & $10^{-4}$ & $10^{-5}$ & $10^{-6}$\\ \midrule\endfirsthead
\toprule $ \Delta t [\text{s}]$ & $10^{-4}$ & $10^{-5}$ & $10^{-6}$ & $10^{-4}$ & $10^{-5}$ & $10^{-6}$ & $10^{-4}$ & $10^{-5}$ & $10^{-6}$\\ \midrule\endhead
\multicolumn{10}{c}{\textbf{\code{(Vm, states, Iion) = (NO, na, NO)}}}\\\hline
$ \Delta x [\text{mm}] $ & \multicolumn{3}{c|}{P8 [ms]} & \multicolumn{3}{c|}{$N_{\mathrm{nl}}$} & \multicolumn{3}{c|}{$N_{\mathrm{lin}}$}\\
0.5 & 55.81 & 59.70 & 59.88 & 3\,000 & 6\,060 & 60\,760 & 12\,000 & 18\,179 & 121\,520\\
0.25 & 45.18 & 50.28 & 50.51 & 2\,231 & 5\,095 & 51\,168 & 15\,615 & 16\,180 & 102\,336\\
0.125 & 40.93 & 46.00 & 46.30 & 1\,662 & 4\,657 & 46\,858 & 18\,281 & 23\,284 & 140\,574\\
\hline
$ \Delta x [\text{mm}] $ & \multicolumn{3}{c|}{wall [min]} & \multicolumn{3}{c|}{CPU [min]} & \multicolumn{3}{c|}{RSS [MiB]}\\
0.5 & 14 & 26 & 61 & 14 & 26 & 61 & 183 & 184 & 183\\
0.25 & 45 & 89 & 231 & 72 & 151 & 855 & 197 & 197 & 198\\
0.125 & 135 & 321 & 782 & 171 & 421 & 1797 & 247 & 247 & 247\\
\hline
\hline
\multicolumn{10}{c}{\textbf{\code{(Vm, states, Iion) = (NO, CCp1, p1)}}}\\\hline
$ \Delta x [\text{mm}] $ & \multicolumn{3}{c|}{P8 [ms]} & \multicolumn{3}{c|}{$N_{\mathrm{nl}}$} & \multicolumn{3}{c|}{$N_{\mathrm{lin}}$}\\
0.5 & 55.81 & 59.70 & 59.88 & 3\,000 & 6\,060 & 60\,760 & 12\,000 & 18\,179 & 121\,520\\
0.25 & 45.18 & 50.28 & 50.51 & 2\,231 & 5\,095 & 51\,168 & 15\,615 & 16\,180 & 102\,336\\
0.125 & 40.93 & 46.00 & 46.30 & 1\,662 & 4\,657 & 46\,858 & 18\,281 & 23\,284 & 140\,574\\
\hline
$ \Delta x [\text{mm}] $ & \multicolumn{3}{c|}{wall [min]} & \multicolumn{3}{c|}{CPU [min]} & \multicolumn{3}{c|}{RSS [MiB]}\\
0.5 & 15 & 26 & 66 & 15 & 26 & 66 & 185 & 186 & 185\\
0.25 & 46 & 90 & 241 & 73 & 152 & 864 & 204 & 204 & 206\\
0.125 & 138 & 330 & 878 & 174 & 429 & 1880 & 273 & 273 & 272\\
\hline
\hline
\multicolumn{10}{c}{\textbf{\code{(Vm, states, Iion) = (p3, CCp1, p1)}}}\\\hline
$ \Delta x [\text{mm}] $ & \multicolumn{3}{c|}{P8 [ms]} & \multicolumn{3}{c|}{$N_{\mathrm{nl}}$} & \multicolumn{3}{c|}{$N_{\mathrm{lin}}$}\\
0.5 & 75.42 & 57.19 & 57.29 & 4\,320 & 7\,956 & 79\,678 & 12\,960 & 23\,868 & 239\,034\\
0.25 & 51.56 & 56.76 & 56.99 & 2\,981 & 5\,741 & 57\,626 & 10\,672 & 17\,223 & 172\,878\\
0.125 & 43.26 & 48.42 & 48.66 & 2\,144 & 4\,904 & 49\,263 & 10\,719 & 14\,712 & 147\,789\\
\hline
$ \Delta x [\text{mm}] $ & \multicolumn{3}{c|}{wall [min]} & \multicolumn{3}{c|}{CPU [min]} & \multicolumn{3}{c|}{RSS [MiB]}\\
0.5 & 23 & 37 & 132 & 75 & 134 & 1097 & 205 & 205 & 205\\
0.25 & 57 & 98 & 299 & 93 & 168 & 999 & 289 & 289 & 291\\
0.125 & 162 & 314 & 978 & 205 & 413 & 2070 & 618 & 619 & 620\\
\hline
\hline
\multicolumn{10}{c}{\textbf{\code{(Vm, states, Iion) = (p3, CCp3, p3)}}}\\\hline
$ \Delta x [\text{mm}] $ & \multicolumn{3}{c|}{P8 [ms]} & \multicolumn{3}{c|}{$N_{\mathrm{nl}}$} & \multicolumn{3}{c|}{$N_{\mathrm{lin}}$}\\
0.5 & 38.88 & 42.27 & 42.55 & 1\,930 & 4\,293 & 43\,188 & 5\,790 & 12\,879 & 129\,564\\
0.25 & 38.93 & 43.30 & 43.77 & 1\,543 & 4\,391 & 44\,368 & 5\,321 & 13\,173 & 133\,104\\
0.125 & 39.80 & 44.71 & 45.30 & 1\,570 & 4\,531 & 45\,883 & 7\,850 & 13\,593 & 137\,649\\
0.1 & 39.86 & 44.81 & 45.40 & 1\,575 & 4\,535 & 45\,925 & 9\,449 & 13\,605 & 137\,775\\
0.05 & 39.62 & 44.61 & -- & 1\,586 & 4\,513 & -- & 15\,856 & 18\,052 & --\\
\hline
$ \Delta x [\text{mm}] $ & \multicolumn{3}{c|}{wall [min]} & \multicolumn{3}{c|}{CPU [min]} & \multicolumn{3}{c|}{RSS [MiB]}\\
0.5 & 15 & 29 & 108 & 38 & 81 & 631 & 219 & 219 & 219\\
0.25 & 41 & 103 & 355 & 59 & 156 & 897 & 331 & 331 & 331\\
0.125 & 154 & 388 & 1409 & 189 & 489 & 2376 & 781 & 793 & 793\\
0.1 & 236 & 578 & 2135 & 271 & 675 & 3122 & 1126 & 1128 & 1126\\
0.05 & 896 & 2216 & -- & 955 & 2341 & -- & 3925 & 3883 & --\\
\hline
\bottomrule\end{longtable}\end{center}
\paragraph{Discussion of the 2D results (\code{Picard}).}
The whole sweep uses \code{Picard}; the \code{diagonalIion} and \code{JFNK} results
in 2D are still pending. Relative to the previous version of this report, the sweep was
extended to $\Delta x\in\{0.1,0.05\}$~mm in the configurations that reach the converged
regime, which for the first time makes it possible to fix a reference value and measure the
error of each configuration against it.

\subparagraph{Converged reference value.}
On the hexahedral mesh, the four paths that reach $\Delta x=\SI{0.05}{mm}$ with
$\Delta t=10^{-6}$ agree on
\begin{equation}
  t_{\mathrm{P8}}^{\ \star} \simeq \SI{45.5}{ms},
  \label{eq:p8ref}
\end{equation}
with a spread of only \SI{0.10}{ms} ($0.2\%$): \SI{45.60}{ms}
(\code{(NO,na,NO)}), \SI{45.50}{ms} (\code{(NO,GP,p1)}), \SI{45.51}{ms}
(\code{(NO,CCp1,p1)} hybrid) and \SI{45.52}{ms} (\code{(p3,CCp3,p3)} hybrid). That
four discretisations with different spatial orders, ionic-source treatments and
state-integration modes converge to the same number is the strongest verification
of the sweep: the solver is consistent, and the differences between
configurations on coarse meshes are discretisation error, not model biases.
All the error percentages that follow are measured against \eqref{eq:p8ref}.

\begin{enumerate}[topsep=3pt,itemsep=3pt,leftmargin=1.5em]
  \item \textbf{High order in the ionic source speeds up and sharpens the front.} Raising
        \code{Iion} from \code{NO} to \code{p3} monotonically lowers the activation
        time: on hexa at $\Delta t=10^{-6}$, $\Delta x=\SI{0.25}{mm}$, P8 goes
        from \SI{67.4}{ms} (\code{(NO,na,NO)}) to \SI{44.4}{ms} (\code{(p3,GP,p3)}). It is
        the same high-order-in-the-source effect as in the previous explicit and implicit
        reports, now resolved with \code{gaussPointODE}.
  \item \textbf{The two families converge from opposite sides, and at very
        different speeds.} With $\Delta t=10^{-5}$~s fixed and refining $\Delta x$ on hexa:
        \begin{center}\footnotesize
        \begin{tabular}{lccccc}
        \toprule
        $\Delta x$ [mm] & 0.5 & 0.25 & 0.125 & 0.1 & 0.05\\
        \midrule
        \code{(NO,na,NO)}            & 157.15 & 67.30 & 50.05 & 47.86 & 45.28\\
        \code{(NO,GP,p1)}            & 94.34  & 58.85 & 48.04 & 46.79 & 45.18\\
        \code{(p3,CCp3,p3)} hybrid   & 36.30  & 44.22 & 45.43 & 45.48 & 45.04\\
        \bottomrule
        \end{tabular}
        \end{center}
        The low-order-source paths approach \eqref{eq:p8ref}
        \emph{from above} and very slowly: \code{(NO,na,NO)} is still $10\%$ high
        at \SI{0.125}{mm} and needs \SI{0.05}{mm} to drop below $1\%$. The high-order
        source path approaches \emph{from below} and is already within
        $0.2\%$ at \SI{0.125}{mm}. In terms of cells: the high-order path
        reaches the target accuracy with $16\times$ fewer cells.
  \item \textbf{$\Delta t=10^{-5}$~s is already enough; $10^{-6}$~s is wasted money.}
        Across all runs that activate, $|P8(10^{-6})-P8(10^{-5})|\le
        \SI{0.5}{ms}$ ($\lesssim 1\%$), whereas the cost is multiplied by
        between $4$ and $10$: on hexa, \code{(p3,CCp3,p3)} hybrid at \SI{0.125}{mm}, the
        wall time goes from $205$ to $845$~min and $N_{\mathrm{nl}}$ from $4\,607$ to $46\,527$,
        to move P8 from \SI{45.43}{} to \SI{45.90}{ms} --- that is, \emph{away}
        from the converged value. By contrast $\Delta t=10^{-4}$~s is \textbf{not}
        enough: it underestimates P8 by $\sim\SI{4.6}{ms}$ ($10\%$) on the high-order
        path.
  \item \textbf{Beware of error cancellation at coarse $\Delta t$.} The
        combination \code{(NO,na,NO)}, $\Delta x=\SI{0.125}{mm}$, $\Delta t=10^{-4}$~s
        gives \SI{45.03}{ms}, apparently a $1\%$ error. It is a coincidence: refining
        $\Delta t$ on the same mesh moves the result to \SI{50.05}{} and
        \SI{50.26}{ms} ($10\%$). The temporal error (which underestimates P8) cancels the
        spatial error (which overestimates it). \textbf{An operating point is only
        valid if it is stable in $\Delta t$}; checking the proximity to the
        converged value is not enough.
  \item \textbf{Pure \code{cellCentredReconstruct} fails when $\Delta t$ is refined.} On
        hexa, the \code{CCp1}/\code{CCp3} configurations without hybrid activate at
        $\Delta t=10^{-4}$~s but \textbf{do not ignite} (\textsc{fail}, $N_a=0$,
        abort by ignition timeout at $t=\SI{0.02}{s}$) at $\Delta t=10^{-5}$ and
        $10^{-6}$~s, and it is now confirmed that the failure persists also at
        $\Delta x=\SI{0.1}{}$ and \SI{0.05}{mm}: \textbf{it is not a problem of
        spatial resolution}. On the triangular mesh the same happens with \code{(p3,CCp3,p3)}.
        The behaviour is governed entirely by $\Delta t$ and is independent of
        $\Delta x$ over a $10\times$ range; its cause is under analysis and is
        not attributed here.
  \item \textbf{The front-aware hybrid rescues \code{CC} at the cost of partial
        \code{GP}.} With \code{frontAwareHybrid true}, the same \code{CCp$k$}
        configurations activate at \emph{all} the $\Delta t$, with activation
        times practically equal to those of \code{gaussPointODE} but much
        cheaper. At $\Delta x=\SI{0.125}{mm}$, $\Delta t=10^{-5}$~s:
        \code{(p3,CCp3,p3)} hybrid gives \SI{45.43}{ms} in $205$~min, against
        \SI{45.40}{ms} in $932$~min for \code{(p3,GP,p3)} without hybrid --- same
        answer, $4.5\times$ cheaper, with the same $N_{\mathrm{nl}}$
        ($\approx 4\,600$). The reduction comes exactly from what was warned above: the
        hybrid integrates the $N_q$ ODEs per cell only in the front band and one
        per cell everywhere else, so that at equal numbers of nonlinear iterations the
        ODE work per evaluation is much smaller. The hybrid therefore combines
        the \emph{correction} of \code{GP} with a good part of the \emph{saving} of
        \code{CC}.
  \item \textbf{Cost scaling.} $N_{\mathrm{nl}}$ and $N_{\mathrm{lin}}$ grow
        $\sim$linearly when refining $\Delta t$ (more steps). The wall time grows
        approximately as $\Delta x^{-2}$ (number of cells in 2D): in
        \code{(p3,CCp3,p3)} hybrid at $\Delta t=10^{-5}$, going from \SI{0.125}{} to
        \SI{0.05}{mm} ($6.25\times$ cells) takes the wall time from $205$ to $1\,123$~min
        ($5.5\times$) and the RSS from $530$ to $2\,367$~MiB ($4.5\times$). The most expensive
        case of the sweep is \code{(p3,GP,p3)} without hybrid on hexa at $\Delta t=10^{-6}$,
        $\Delta x=\SI{0.125}{mm}$: $2\,228$~min ($\sim 37$~h). The RSS grows with the
        LRE stencil and with $1/\Delta x^2$, and is essentially independent of
        $\Delta t$.
  \item \textbf{Wall vs.\ CPU: why the deterministic quantities are reported.} At
        $\Delta x\le\SI{0.125}{mm}$ the CPU time exceeds the wall time (multi-threaded
        execution under contention; e.g.\ \code{(NO,na,NO)} hexa $\Delta t=10^{-6}$:
        wall $339$~min, CPU $1098$~min), whereas on coarse meshes they coincide. That
        divergence is precisely the machine-dependent noise that motivates using
        $N_{\mathrm{nl}}$ and $N_{\mathrm{lin}}$ as a reproducible cost reference.
  \item \textbf{Triangular mesh.} The qualitative pattern is preserved (failure of
        pure \code{CC}, rescue by the hybrid, high order in $\Iion$ speeding things up).
        Three observations of its own: (i) on the triangular mesh, the configurations with
        \code{Iion}$=$\code{p1} coincide \emph{exactly} with the baseline
        \code{(NO,na,NO)} (same P8, same $N_{\mathrm{nl}}$, same
        $N_{\mathrm{lin}}$) --- the $p1$ reconstruction of $\Iion$ contributes nothing on this
        unstructured mesh, and only \code{p3} separates; (ii) the converged
        triangular value, $\simeq\SI{44.6}{ms}$ for \code{(p3,CCp3,p3)} hybrid at
        \SI{0.05}{mm}, lies $\sim 2\%$ below the hexahedral one
        \eqref{eq:p8ref} --- a genuine mesh difference, not a method one; (iii) the
        larger number of cells per $\Delta x$ (5\,954 vs 2\,560 at \SI{0.25}{mm})
        proportionally increases every quantity, so that at equal
        $\Delta x$ the triangular case costs between $1.9$ and $2.4\times$ more wall time.
\end{enumerate}

\subparagraph{Recommended configuration.}
\label{sec:res2d:reco}
Combining accuracy and cost, and taking \eqref{eq:p8ref} as the reference, the sweep
identifies one clearly dominant operating point and two alternatives depending on the
budget:

\begin{center}\footnotesize
\begin{tabular}{llccrrr}
\toprule
Use & Triplet / hybrid & $\Delta x$ [mm] & $\Delta t$ [s] & P8 [ms] & error & wall [min]\\
\midrule
\textbf{Recommended} & \code{(p3,CCp3,p3)} + hybrid & 0.125 & $10^{-5}$ & 45.43 & $0.2\%$ & \textbf{205}\\
Quick exploration    & \code{(p3,CCp3,p3)} + hybrid & 0.25  & $10^{-5}$ & 44.22 & $2.8\%$ & \textbf{60}\\
Reference            & \code{(p3,CCp3,p3)} + hybrid & 0.05  & $10^{-6}$ & 45.52 & $0.05\%$ & 4\,211\\
\midrule
\multicolumn{7}{l}{\emph{Alternatives at the same error level, for comparison:}}\\
Low order            & \code{(NO,na,NO)}            & 0.05  & $10^{-5}$ & 45.28 & $0.5\%$ & 739\\
High order, no hybrid & \code{(p3,GP,p3)}           & 0.125 & $10^{-5}$ & 45.40 & $0.2\%$ & 932\\
\bottomrule
\end{tabular}
\end{center}

\noindent\textbf{Why that combination.} The three ingredients are all necessary and each
one solves a different problem:
\begin{itemize}[topsep=2pt,itemsep=2pt,leftmargin=1.4em]
  \item \code{Iion}$=$\code{p3} with \code{CCp3} states gives the fast spatial
        convergence (point 2): it is what allows stopping at \SI{0.125}{mm} instead of
        \SI{0.05}{mm}.
  \item \code{frontAwareHybrid true} is what makes \code{CCp3} \emph{ignite}
        (point 5) and, at the same time, avoids paying the $N_q$ ODEs per cell over the whole
        domain (point 6).
  \item $\Delta t=10^{-5}$~s is the coarsest step that gives temporal stability without
        overspending (point 3).
\end{itemize}

\noindent\textbf{Gain.} Against the fully resolved reference, the recommended
configuration is $20\times$ cheaper ($205$ against $4\,211$~min) for $0.2\%$ error.
Against the low-order alternative at equal accuracy it is $3.6\times$ cheaper ($205$
against $739$~min) and uses $16\times$ fewer cells; against the same physics without hybrid,
$4.5\times$ cheaper. If only a qualitative sweep is needed, dropping to
\SI{0.25}{mm} costs $60$~min ($70\times$ less than the reference) in exchange for an error
of $2.8\%$, which is still smaller than the $10\%$ of the classical baseline at \SI{0.125}{mm}.

\noindent\textbf{What to avoid.} (i) \code{CCp$k$} without hybrid at any
$\Delta t\le 10^{-5}$~s: it does not ignite, and the failure is not fixed by refining the mesh.
(ii) $\Delta t=10^{-4}$~s as a production operating point: it can give an apparently
good number by error cancellation (point 4) without being reproducible under refinement.
(iii) $\Delta t=10^{-6}$~s: it multiplies the cost by $4$--$10$ without improving the answer.
(iv) \code{(p3,GP,p3)} without hybrid: it gives the same answer as the recommended one but
$4.5\times$ more expensive.
\subsection{3D Niederer-style benchmark}
\label{sec:results:3D}

The 3D sweep follows Niederer's published slab geometry
($20\times 3\times 7$~mm, $40\times 6\times 14$ cells at $\Delta x = 0.5$~mm). At the
time of writing, only the \code{diagonalIion} and \code{Picard} sweeps include
3D runs, all at $\Delta x = 0.5$~mm. The 3D \code{JFNK} sweep has not been run yet.

\begin{center}\footnotesize
\begin{longtable}{llrr}
\caption{3D P8 activation time [ms] at $\Delta x = 0.5$~mm,
$\Delta t = 10^{-4}$~s. The columns are the available nonlinear methods.}\\
\toprule
Vm & Iion & \code{Picard} & \code{diagonalIion} \\
\midrule\endfirsthead\toprule
Vm & Iion & \code{Picard} & \code{diagonalIion} \\
\midrule\endhead
NO & NO & 136.67 & 136.61 \\
NO & p1 &  33.50 &  33.39 \\
NO & p3 &  30.59 &  30.51 \\
p3 & NO & 149.13 & 149.06 \\
p3 & p1 &  32.44 &  32.36 \\
p3 & p3 &  29.63 & NC     \\
\bottomrule\end{longtable}\end{center}

\begin{center}\footnotesize
\begin{longtable}{llrr}
\caption{3D solver wall-clock time [min] at $\Delta x = 0.5$~mm,
$\Delta t = 10^{-4}$~s.}\\
\toprule
Vm & Iion & \code{Picard} & \code{diagonalIion} \\
\midrule\endfirsthead\toprule
Vm & Iion & \code{Picard} & \code{diagonalIion} \\
\midrule\endhead
NO & NO &   80 &   82 \\
NO & p1 &  168 &  229 \\
NO & p3 & 1489 & 2159 \\
p3 & NO &  163 &  143 \\
p3 & p1 &  167 &  269 \\
p3 & p3 & 1487 & NC   \\
\bottomrule\end{longtable}\end{center}

\begin{center}\footnotesize
\begin{longtable}{llrr}
\caption{3D peak resident set size [MiB] at $\Delta x = 0.5$~mm,
$\Delta t = 10^{-4}$~s.}\\
\toprule
Vm & Iion & \code{Picard} & \code{diagonalIion} \\
\midrule\endfirsthead\toprule
Vm & Iion & \code{Picard} & \code{diagonalIion} \\
\midrule\endhead
NO & NO &  156 &  157 \\
NO & p1 &  211 &  211 \\
NO & p3 &  787 &  849 \\
p3 & NO & 1123 & 1121 \\
p3 & p1 & 1164 & 1162 \\
p3 & p3 & 1565 & NC   \\
\bottomrule\end{longtable}\end{center}
\paragraph{Discussion of the 3D results.}
On the coarse 3D slab ($\Delta x = 0.5$~mm), enabling LRE only in the ionic source
changes the activation time from $\approx 136.7$~ms (\code{NO}/\code{NO}) to
$\approx 33.5$~ms (\code{NO}/\code{p1}) and $\approx 30.6$~ms (\code{NO}/\code{p3}) ---
the same qualitative high-order-in-the-source effect already seen in the 2D sweep and in
the explicit/implicit Niederer reports. The fully high-order case
(\code{p3}/\code{p3}) reaches $29.63$~ms with Picard (the diagonalIion counterpart has
not been completed yet). The two nonlinear methods again agree to better than
$0.1$~ms for every entry where both are available; Picard is generally
faster, with the gap most noticeable in the heaviest case (\code{NO}/\code{p3}:
$1489$~min for Picard vs.\ $2159$~min for diagonalIion). The peak RSS scales with the
LRE stencil size and is essentially independent of the method. The data set
is too sparse in $\Delta x$ to draw quantitative conclusions about mesh
convergence in 3D; the immediate gap to fill is the missing
\code{p3}/\code{p3} diagonalIion entry and a 3D JFNK sweep on the same mesh.

% ================================================================== %
\appendix
\section{Dimension table}
% ================================================================== %

\begin{center}
\begin{tabular}{@{}lll@{}}
\toprule
Field & OpenFOAM \code{dimensionSet} & SI \\
\midrule
\code{Vm} & $[1, 2, -3, 0, 0, -1, 0]$ & volt \\
\code{Iion} & \code{dimVoltage/dimTime} & \si{\volt\per\second} \\
\code{externalStimulusCurrent} & \code{dimCurrent/dimVolume} &
    \si{\ampere\per\meter\cubed} \\
\code{conductivity} &
    \code{dimCurrent\^{}2 dimTime\^{}3 / (dimMass dimVolume)} &
    \si{\siemens\per\meter} \\
\code{chi} & \code{dimless/dimLength} & \si{\per\meter} \\
\code{Cm} & \code{dimCurrent dimTime / (dimVoltage dimArea)} &
    \si{\farad\per\meter\squared} \\
\bottomrule
\end{tabular}
\end{center}

\section{File map}

\begin{center}
\begin{tabular}{@{}p{0.55\linewidth}p{0.40\linewidth}@{}}
\toprule
Path (relative to \code{cardiacFoam-pc/}) & Role \\
\midrule
\code{applications/solvers/highOrderElectroActivationFoamImplicitPETScDistributed/}
   \code{highOrderElectroActivationFoamImplicitPETScDistributed.C}
   & Main translation unit (time loop, nonlinear solvers, Eigen/PETSc linear
   back-ends, distributed \code{MPIAIJ} operators, LRE reconstruction,
   activation tracking). \\
\code{applications/solvers/highOrderElectroActivationFoamImplicitPETSc/}
   & Serial ancestor, and \code{...PETScParallel} the intermediate step
   (parallel stencils, serial algebra). Kept as references
   (Section~\ref{sec:distributed}). \\
\code{applications/solvers/highOrderManufacturedFDAImplicitPETScDistributed/}
   & Manufactured-solution sibling sharing the parallel design; it carries the
   coupled-face stabilisation exchange and the \code{CF\_DUMP\_K} /
   \code{CF\_DUMP\_STENCILS} diagnostics (Section~\ref{sec:dist-verif}). \\
\code{src/highOrderAdapter/highOrderInterp.H} & Adapter over solids4foam's
   \code{movingLeastSquares}: stencils with global columns, halo exchange,
   reconstruction coefficients. \\
\code{src/ionicModels/} & Ionic models (TNNP and others): embedded CellML kernel,
   state update, current evaluation and local ODE integration. \\
\code{applications/solvers/.../createFields.H} & Field initialisation
   and dictionary reading. \\
\code{tutorials\_electro/.../run\_convergence.py} & Driver script. \\
\code{constant/cardiacProperties} & $\chi$, $C_{m}$, $\Dten$,
   ionic model. \\
\code{constant/timeIntegrationProperties} & ODE solver, tolerances,
   sub-step cap. \\
\code{constant/spatialIntegrationProperties} & Time scheme, HO flags,
   JFNK / linear / line-search options. \\
\code{constant/stimulusProtocol} & Stimulus regions and timing. \\
\bottomrule
\end{tabular}
\end{center}

% ================================================================== %
\section{Linear-algebra back-ends and sibling solver}
\label{sec:backends}
% ================================================================== %

This solver \emph{integrates both linear-algebra back-ends} in a single
executable, selectable at run time with \code{linearSolverBackend} (assembled
branches) and \code{jfnkLinearSolverBackend} (JFNK Krylov). Matrix assembly and
storage always use \code{Eigen::SparseMatrix}; the back-end only
determines how the system is solved.

\begin{center}
\begin{tabular}{@{}p{0.34\linewidth}p{0.30\linewidth}p{0.30\linewidth}@{}}
\toprule
Aspect & \code{...Backend = Eigen} & \code{...Backend = PETSc} \\
\midrule
Matrix container & \code{Eigen::SparseMatrix} & idem (copied to \code{Mat}) \\
Sparse direct solver & \code{Eigen::SparseLU} & \code{KSPPREONLY+PCLU} \\
Iterative solver (Picard/diag) & \code{Eigen::BiCGSTAB+ILUT} &
    \code{KSP+PC} (\code{gmres}+\code{ilu}, \code{cg}+\code{gamg}, \dots) \\
JFNK inner Krylov & In-house GMRES & \code{KSPGMRES} over \code{MatShell} \\
JFNK preconditioner & \multicolumn{2}{c}{\code{jfnkPreconditioner}: \code{none} or
    \code{diffusion} (ILU of $\Aimp$) --- available in both} \\
\bottomrule
\end{tabular}
\end{center}

The convergence parameters (\code{nonlinearTolerance}, \code{jfnkEpsilon},
\code{jfnkLineSearch*}, \code{jfnkInitGuess*}, \code{jfnkClampODEInput},
\code{nonlinearAcceptUnconverged}, \code{ODE\_*}) have the same meaning
regardless of the back-end, so that the same \code{run\_convergence.py}
configuration produces an algorithmically equivalent run (apart from
the performance differences of the linear solve).

\paragraph{Sibling solver \code{highOrderElectroActivationFoamImplicit}.}
There is also a simpler version, with an \emph{Eigen-only} back-end and without the
optimisations of Section~\ref{sec:optim}. Its main algorithmic difference is
that it runs a \emph{fixed} number of nonlinear iterations (\code{nonlinearIterations},
$1$ by default) \emph{without} a convergence test: with the default value it is a
semi-implicit scheme with \emph{one} ODE evaluation and \emph{one} linear solve per
time step. This makes it much faster at the cost of not controlling the nonlinear
residual; at the small time steps of the benchmark the difference in the results
is marginal (see Sections~\ref{sec:optim} and~\ref{sec:exp}).

% ================================================================== %
\section{Distributed-memory parallelisation}
\label{sec:distributed}
% ================================================================== %

This variant, \code{highOrderElectroActivationFoamImplicitPETScDistributed},
leaves every algorithm of the previous sections unchanged and replaces the
linear algebra with a genuinely distributed one. The distinction matters,
because there are three solvers in the family and only the last one is parallel
in the algebraic sense:

\begin{center}
\begin{tabular}{@{}p{0.30\linewidth}p{0.62\linewidth}@{}}
\toprule
Variant & Linear algebra \\
\midrule
\code{...ImplicitPETSc} & Serial mesh, serial LRE, Eigen assembly, PETSc as a
   solver on \code{PETSC\_COMM\_SELF}. \\
\code{...ImplicitPETScParallel} & Parallel high-order library: the mesh is
   decomposed and the stencils cross ranks, but the assembly still produces one
   Eigen matrix per rank and hands the finished CSR to PETSc on
   \code{PETSC\_COMM\_SELF}. Serial by construction. \\
\code{...ImplicitPETScDistributed} & Matrices are \code{MPIAIJ} on
   \code{PETSC\_COMM\_WORLD} and every matrix--vector product goes through
   \code{MatMult}. \\
\bottomrule
\end{tabular}
\end{center}

The three are kept side by side deliberately: each is the reference against
which the next one is checked, rather than being compared from memory.

% ------------------------------------------------------------------ %
\subsection{Matrix layout: local rows, global columns}
\label{sec:dist-layout}
% ------------------------------------------------------------------ %

Each rank owns the rows of its own cells and nothing else. Columns, in contrast,
are \emph{global} cell indices, taken from the \code{globalIndex} that the
high-order library builds together with its stencils. Writing
$n_{\text{loc}}$ for this rank's cell count and $n_{\text{glob}}$ for the total,
every assembled operator has shape $n_{\text{loc}}\times n_{\text{glob}}$, and
the row offset of this rank is
\[
   \texttt{gRowStart} \;=\; \texttt{globalIndex::localStart()} .
\]

That is exactly PETSc's \code{MPIAIJ} layout, so the translation from the Eigen
container to the distributed matrix (\code{buildDistributedMat}) only has to
offset the row index and split the preallocation into the diagonal block
(columns this rank owns) and the off-diagonal one. Getting that split right is
what stops assembly from degenerating into repeated reallocation.

In serial $\texttt{gRowStart}=0$ and $n_{\text{glob}}=n_{\text{loc}}$, so the
same code path produces exactly the sequential matrix it replaces.

\paragraph{A bug family worth naming.} That last property is convenient and
dangerous in equal measure: \emph{a local index used where a global one is
expected is invisible in serial}. Five instances were found in this family of
solvers, all of them silent at one rank:

\begin{center}
\begin{tabular}{@{}p{0.42\linewidth}p{0.50\linewidth}@{}}
\toprule
Site & Symptom at $>1$ rank \\
\midrule
Diagonal of the mass matrix & Extra nonzeros in $\Aimp$, exactly the per-rank
   row count. \\
\code{addDiagonalToMatrix} column & As above. \\
\code{updateValues} size check & Accepts a vector of the wrong length. \\
\code{updateValues} row offset & Values written into the wrong rows. \\
Orthogonal assembly of the manufactured sibling & Local columns throughout,
   so the operator could not be assembled correctly at all. \\
\bottomrule
\end{tabular}
\end{center}

Section~\ref{sec:dist-verif} explains why the usual sanity checks do not see
any of these.

% ------------------------------------------------------------------ %
\subsection{What has to cross a processor boundary}
\label{sec:dist-exchange}
% ------------------------------------------------------------------ %

Three kinds of data cross a partition, and each needs a different mechanism.

\paragraph{(i) Reconstruction stencils --- nothing to do.}
The high-order library builds its stencils geometrically over a halo whose depth
is controlled by \dictkey{haloDepthScale}, and emits \emph{global} column
indices. The flux term at a processor face therefore assembles exactly as at an
internal face: the stencil already reaches across the cut.

\paragraph{(ii) Cell-centred data --- \code{syncTools}.}
The conductivity of the cell on the far side of a coupled face, and its global
index, are obtained with \code{syncTools::swapBoundaryCellList}, indexed by
boundary-face position. Two tempting alternatives are both wrong:

\begin{itemize}
  \item a processor patch \emph{field} holds the interpolated value \emph{at the
        face}, not the value of the neighbouring \emph{cell};
  \item \code{mesh.C()} is a \code{slicedVolVectorField} whose coupled values are
        only valid after an evaluation that the assembly does not perform.
\end{itemize}

For the same reason the owner-to-neighbour vector $\mathbf{d}_{PN}$ is taken
from \code{patch.delta()}, which for a \code{processorFvPatch} returns
$\mathbf{C}_{N}-\mathbf{C}_{P}$ built from both sides' geometry, and never from
\code{C[nei] - C[own]}.

\paragraph{(iii) Variable-length matrix rows --- \code{PstreamBuffers}.}
The stabilisation jump is the one term that needs the neighbouring cell's whole
reconstruction row: up to $\sim 60$ (global column, coefficient) pairs for a
$p3$ stencil in 3-D. No field swap can carry that.

The mechanism, implemented in the manufactured sibling
\code{highOrderManufacturedFDAImplicitPETScDistributed}, is to let the rank that
\emph{owns} the cell build the row itself and ship the pairs, evaluated at the
face centre --- a point both ranks have. Nothing else is exchanged: no cell
centres, no stencils, no coefficient tables. Writing the jump from the point of
view of the local cell,
\[
   \frac{a}{V_{c}}\Bigl( T_{\text{far}}(\mathbf{x}_{f})
                        - T_{c}(\mathbf{x}_{f}) \Bigr),
\]
where $T_{c}$ is the Taylor reconstruction of cell $c$, also removes the
owner/neighbour convention: the $+$ that the internal-face loop writes into the
owner row and the $-$ it writes into the neighbour row are the same expression
seen from either cell, and the face coefficient $a$ depends on the orientation
only through $|\mathbf{d}_{PN}\!\cdot\!\mathbf{n}|$ and
$\mathbf{n}\!\cdot\!\Dten\!\cdot\!\mathbf{n}$. Both ranks therefore compute the
same $a$ without agreeing on anything.

\paragraph{Status in this solver.} The coupled branch of the orthogonal assembly
adds the flux term and \emph{not} the Rhie--Chow jump. This is inert at the
default $\alpha=0$ that the benchmarks use, and it is documented in the source
rather than silently skipped; porting the exchange described above is the
outstanding item.

% ------------------------------------------------------------------ %
\subsection{Global reductions}
% ------------------------------------------------------------------ %

Everything that used to be a loop over all cells and is now a loop over this
rank's cells needs a reduction: the error norms of
Section~\ref{sec:results} (\code{gSum} over volumes and squared errors), the
nonlinear residuals and their stopping tests, and the activation-time sampling.

The one that is easy to miss is the \emph{Armijo backtracking line search},
condition~\eqref{eq:armijo}. Its accepted step must be a single number for the
whole domain; without the reductions inside the search, each rank would accept a
different $\lambda$ and the Newton iterate would stop being a function of the
mesh alone. That the parallel run reproduces the serial \code{lineSearchIters}
\emph{exactly} is what validates this.

% ------------------------------------------------------------------ %
\subsection{Verifying a parallel operator}
\label{sec:dist-verif}
% ------------------------------------------------------------------ %

Two checks are cheap and are \emph{not} sufficient on their own:

\begin{itemize}
  \item \textbf{Row sums.} With no-flux boundaries a constant field has zero
        gradient, so every row of $\Kstif$ must sum to zero. This catches an
        empty or malformed operator --- it caught a
        \code{CompactListList::operator[]} returning a view by value, which
        compiled but produced an empty stencil --- but it cannot see an
        \emph{omitted} face, because omitting a face contributes nothing to
        either of its rows and the sums still vanish.
  \item \textbf{Operator fingerprint} (global nnz, $\sum a_{ij}$,
        $\sum|a_{ij}|$, $\max_i |\sum_j a_{ij}|$). Since the stencils are
        geometric, the global nnz must match the serial value exactly. But a
        value written into the \emph{wrong column} changes neither nnz, nor any
        sum, nor any row sum.
\end{itemize}

The check that does decide is an entry-by-entry comparison. With
\code{CF\_DUMP\_K} set in the environment, the manufactured sibling writes
$\Kstif$ as (global row, global column, value), one file per rank; mapping the
parallel indices back through
\code{processor$N$/constant/polyMesh/cellProcAddressing} gives a direct diff
against the serial assembly. \code{CF\_DUMP\_STENCILS} does the same one level
lower, for stencil \emph{membership} and cell centres, which is what
distinguishes a differing coefficient from a differing stencil.

\paragraph{Comparing iteration by iteration.} Use
\dictkey{petscLinearPcType} \code{jacobi}. PETSc wraps \code{ilu} and \code{lu}
in block-Jacobi in parallel, so those preconditioners are partition-dependent
\emph{by construction}: with \code{ilu} a converged manufactured case shows a
$4.5\times10^{-6}$ deviation between one and two ranks that is entirely
algebraic, and with \code{jacobi} the same case agrees to all printed digits.
Mistaking that for a discretisation error costs a lot of time.

% ------------------------------------------------------------------ %
\subsection{Partition dependence of the MLS stencils}
\label{sec:dist-stencils}
% ------------------------------------------------------------------ %

One genuine parallel inconsistency remains, and it lives in the stencil
selection of the high-order library, not in any solver.

\paragraph{Mechanism.} A stencil is meant to be \emph{the $N$ cells whose
centres are nearest the face centre}. The library reaches that in two steps: it
first builds a pool of candidates by growing complete layers of
\code{cellCells} until the pool holds $\lceil 1.5N \rceil$ cells, and then keeps
the nearest $N$ of the pool by Euclidean distance. The second step is geometric;
\textbf{the first is topological}. On a tetrahedral mesh ``within $k$ hops'' and
``within radius $r$'' are not the same set, and with only $0.5N$ of slack the
last positions of the stencil are decided by the connectivity rather than by the
geometry.

Under decomposition the layer growth stops at the cut, and the far side is
supplied instead by a query that returns every cell of the neighbouring rank
inside a bounding box --- a \emph{geometric} set, and a generous one. Near a
partition the selection therefore mixes a topological candidate set on one side
with a geometric one on the other. This is why the deviation does not respond to
\dictkey{haloDepthScale}: the box was never the binding constraint.

\paragraph{Measurements} (3-D unstructured tetrahedra, 4971 cells, one rank
against two, $\alpha=0$ so that no stabilisation is involved):

\begin{center}
\begin{tabular}{@{}p{0.55\linewidth}p{0.37\linewidth}@{}}
\toprule
Observable & Value \\
\midrule
Face stencils that differ & 536 of 10676 \\
Their size & identical (80): members are \emph{substituted}, not lost \\
Position of the swapped members & 98\,\% in the last 10\,\% of the
   distance-ordered stencil \\
Rows of $\Kstif$ that differ & 552 of 4971, up to 4.17 against
   $\max|K|=256$ \\
Cut faces whose two ranks disagree with \emph{each other} & 9 \\
Cell centres, serial vs decomposed & up to $3.3\times10^{-16}$ (the centroid sum
   runs over a renumbered face list) \\
\bottomrule
\end{tabular}
\end{center}

The nine cut faces are a defect on their own terms, with no serial comparison
needed: the two rows of one face are then built from different flux stencils, so
the flux leaving a cell is not the flux entering its neighbour and the scheme is
not conservative there.

\paragraph{How much it matters.} Comparing the $V_{\mathrm{m}}$ \emph{fields}
cell by cell rather than the deviation of a norm, at $\alpha=0.1$ with
\code{jacobi}:

\[
  \frac{\norm{V_{\mathrm{m}}^{np=1}-V_{\mathrm{m}}^{np=2}}}
       {\norm{V_{\mathrm{m}}-V_{\mathrm{m}}^{\text{exact}}}}
  = \frac{1.51\times10^{-5}}{6.84\times10^{-4}} = 2.2\,\%,
  \qquad
  \frac{\norm{\Delta V_{\mathrm{m}}}}{\text{amplitude}} = 7.5\times10^{-6}.
\]

Over a refinement sweep $N=8,10,12,14$ the deviation stays at the
$10^{-4}$--$10^{-3}$ level while the error itself falls from
$1.19\times10^{-3}$ to $2.28\times10^{-4}$: the perturbation shrinks at the same
rate as the error. Consequently the fitted convergence order is untouched,

\[
  p^{np=1}_{V_{\mathrm{m}},L_2} = 2.90722821,
  \qquad
  p^{np=2}_{V_{\mathrm{m}},L_2} = 2.90722635,
\]

a difference of $1.9\times10^{-6}$ against an acceptance tolerance of $0.05$.
That is what two \emph{consistent} schemes of the same formal order predict:
their difference is a same-order perturbation, not a systematic error.

\paragraph{How it would be fixed.}
\begin{enumerate}
  \item Make the candidate pool geometric on both sides --- query cell centres
        by radius (\code{indexedOctree}) instead of growing \code{cellCells}
        layers, with the radius taken from the layer estimate as an upper bound.
        The stencil then becomes a function of (face centre, cell centres, $N$)
        alone, invariant under partitioning.
  \item Stop the tie clipping: the parallel branch filters candidates by radius
        \emph{before} the equal-distance expansion runs, so a tie group
        straddling that radius is cut in half (17 members in serial against 16 at
        two ranks, on a structured 2-D case). Inflating the radius by
        $(1+\texttt{relTol})$ and routing the serial early return through the
        same selection code closes it.
  \item Assert completeness instead of failing silently: check that the $N$-th
        distance lies strictly inside the queried radius.
  \item One face, one stencil: have the lower-ranked side of a cut face own the
        stencil and send it, which restores conservation across the cut.
\end{enumerate}

\code{relTol} ($10^{-8}$), \code{overSampleFactor} ($1.5$) and
\code{searchExpansionFactor} ($1.3$) are hard-coded default arguments, not
dictionary entries, so none of this can be explored from a case: it is a library
change.

\paragraph{A caveat before touching it.} The serial run is the truncated one ---
its candidate pool is purely topological everywhere, not only near a cut --- so
fixing the selection changes serial results too, and with them every convergence
table produced so far. That is a decision to take deliberately, not a side
effect to absorb.

% ================================================================== %

\subsection{Verification cases: what was run and where it lives}
\label{sec:elec-parverif-cases}

Every case below sits under
\code{tutorials\_electro/electroTestParallelVerification/}. The driver is
\code{tutorials\_electro/electroDistValidation/}; the directories here are
copies of it with one setting changed each.

\begin{center}
\begin{tabular}{@{}llp{6.2cm}@{}}
\toprule
Directory & Test & What it establishes \\
\midrule
\code{electroA3\_np1}, \code{\_np2}, \code{\_np4}
  & A3, physical acceptance
  & activation time at P8 and conduction velocity, full run to activation \\
\code{electroA4\_np2}, \code{\_np4}
  & A4, decomposition invariance
  & \code{simple} instead of \code{scotch}: a different cut geometry \\
\code{electroTri\_np1}, \code{\_np2}, \code{\_np4}
  & unstructured mesh
  & the gmsh route, 1500 cells \\
\code{matrix\_el\_2D\_np1}, \code{\_np2}, \code{matrix\_el\_3D\_*}
  & method matrix
  & Picard / JFNK / diagonalIion, hexa + triangular, 2-D and 3-D \\
\code{jac\_2D\_np1}, \code{jac\_2D\_np2}
  & preconditioner isolation
  & \code{jacobi}, the only rank-independent preconditioner \\
\code{perf\_\{ilu,jac,gamg,hypre\}\_np\{1,2\}}
  & linear-solver performance
  & $\Delta x = 0.5$\,mm, 640 cells \\
\code{perf025\_\{ilu,jac,gamg,hypre\}\_np\{1,2\}}
  & idem, refined
  & $\Delta x = 0.25$\,mm, 2560 cells \\
\code{pre\_\{0p0,0p1\}\_np\{1,2\}}
  & stabilisation, before
  & built with the \code{electroPreStab} binary \\
\code{post\_\{0p0,0p1\}\_np\{1,2\}}
  & stabilisation, after
  & built with the ported binary \\
\code{stab\_\{0p0,0p1\}\_np\{1,2\}}
  & first baseline attempt
  & \emph{discarded}: rebuilds landed mid-sweep \\
\bottomrule
\end{tabular}
\end{center}

\paragraph{A3, physical acceptance} (\code{electroA3\_*}). Hexa,
$\Delta x = 0.5$\,mm, $\Delta t = 10^{-4}$, triad $(p3,\mathrm{CCp3},p3)$, run
until the wavefront reaches P8. Sequential runs on an idle machine, so the wall
times are valid.

\begin{center}
\begin{tabular}{@{}rlrrrrr@{}}
\toprule
Ranks & Status & $T_{P8}$ [ms] & dev. & CV [m/s] & Wall [s] & Speed-up \\
\midrule
1 & OK & 52.0795930112 & ---            & 0.38388 & 310.3 & $1.00\times$ \\
2 & OK & 52.0795930123 & \num{2.11e-11} & 0.38388 & 177.0 & $1.75\times$ \\
4 & OK & 52.0795930133 & \num{4.03e-11} & 0.38388 & 102.9 & $3.01\times$ \\
\bottomrule
\end{tabular}
\end{center}

\code{OK} is the most important entry in each row: it is the binary
conduction-block detector, and it says the wavefront crosses every partition
cut. The $0.5\,\%$ tolerance proposed for $|\Delta T_{P8}|/T_{P8}$ was
\emph{wrong by eight orders of magnitude} --- the observed spread is
\num{4e-11}. Picard converges to the same fixed point deterministically, so a
loose stopping residual introduces no drift.

\paragraph{A4, decomposition invariance} (\code{electroA4\_*}). The same run cut
with \code{simple} ($n = (2\,1\,1)$ and $(2\,2\,1)$) rather than \code{scotch}:
\code{OK} in both, $T_{P8}$ within \num{9.22e-11} of serial, CV identical to
five decimals. Two unrelated cut geometries agreeing is far stronger evidence
than several rank counts under one method.

\paragraph{Unstructured mesh} (\code{electroTri\_*}). 1500 cells through the
gmsh route: \code{OK} at 1, 2 and 4 ranks, $T_{P8}$ within \num{6.55e-11},
speed-up $1.65\times$ and $2.84\times$.

\paragraph{Linear-solver performance} (\code{perf\_*}, \code{perf025\_*}).
Krylov iterations per nonlinear iteration vary by more than a factor three
between preconditioners while the wall times agree to within 8\,\%, at both mesh
sizes. \emph{The linear solve is not the bottleneck} --- the ionic ODEs are ---
so changing preconditioner buys nothing here. Refining $0.5 \to 0.25$\,mm
multiplies the cost by $4.3$, i.e. essentially linearly in cell count, and the
iteration counts stay flat (ilu 3.0, hypre 4.0, jacobi 8.0), so no crossover
appears. \code{gamg} failed on the finer mesh and was not diagnosed.

\paragraph{Preconditioner isolation} (\code{jac\_2D\_*}). \code{ilu} has no
MPIAIJ implementation and is wrapped in block-Jacobi in parallel, which is
partition-dependent by construction; \code{jacobi} is purely diagonal and
therefore the same operator at every rank count. Switching to it collapsed the
$np=1$ vs $np=2$ trajectory differences: five of six configurations became
\emph{identical}, and the row counts matched in all six (\code{triangular} with
Picard went from $184$ vs $275$ rows to $184$ vs $184$).

\paragraph{Processor-face stabilisation} (\code{pre\_*} vs \code{post\_*}).
Deviation of the residual trace between $np=1$ and $np=2$ at $\alpha=0.1$:

\begin{center}
\begin{tabular}{@{}llrrrr@{}}
\toprule
Mesh & Method & Before & After & $\alpha=0$ control & Gain \\
\midrule
hexa       & Picard       & \num{3.23e-09} & \num{9.06e-12} & \num{8.53e-12} & $356\times$ \\
hexa       & JFNK         & \num{2.58e-05} & \num{2.59e-05} & \num{2.48e-05} & $1\times$ \\
hexa       & diagonalIion & \num{1.27e-08} & \num{4.51e-10} & \num{2.28e-10} & $28\times$ \\
triangular & Picard       & \num{2.45e-07} & \num{4.68e-11} & \num{9.55e-12} & $5240\times$ \\
triangular & JFNK         & \num{4.44e-06} & \num{4.43e-06} & \num{4.64e-06} & $1\times$ \\
triangular & diagonalIion & \num{1.59e-07} & \num{6.17e-09} & \num{5.59e-10} & $26\times$ \\
\bottomrule
\end{tabular}
\end{center}

Serial output is bit-identical (SHA-256) in all twelve configurations, at both
$\alpha$ values. The assembled paths collapse onto the $\alpha=0$ control, which
is the achievable floor. JFNK does not move, and that is expected: its deviation
is dominated by its own preconditioner --- the Eigen ILUT of the diagonal block,
partition-dependent by construction --- not by the assembled stabilisation. The
same \num{2.48e-05} appears with \code{ilu} and with \code{jacobi}, which is what
identified it.

Two \emph{distinct} binaries were built for this comparison,
\code{electroPreStab} and the ported solver, because a first attempt
(\code{stab\_*}) had rebuilds landing while the sweep was still running. That
matters: had the ``before'' already contained the fix, the ``after'' would have
been identical, and a null result is indistinguishable from a working one.


\section{Performance optimisations}
\label{sec:optim}
% ================================================================== %

This section documents four optimisations added to the solver. All of them are
\emph{additive} and are governed by dictionary keys whose default
value reproduces \textbf{exactly} the previous behaviour: no existing
run changes its result unless they are explicitly enabled. They also update
some rows of the comparison table of the previous Section (the assembled
back-end now reuses the factorisation, JFNK admits a preconditioner,
and the ionic ODE can be integrated in parallel with OpenMP).

\subsection{Matrix and factorisation reuse (Phase B)}
\label{sec:optim-B}

The implicit operator $\Aimp = \Mmass/\Delta t - \theta\,\Lop$ depends only on the
constant matrices $\Mmass$, $\Lop$, on the parameter $\theta$ and on the step $\Delta t$.
Since $\Delta t$ is fixed except in the last clipped step, $\Aimp$ is constant
during almost the whole simulation. Previously, $\Aimp$ was reassembled and
\textbf{re-factorised} at every nonlinear iteration; now it is assembled once and the
factorisation (PETSc KSP, or \code{Eigen::SparseLU}/\code{BiCGSTAB}) is kept
for the whole run. An update policy distinguishes three cases per solve:
\begin{itemize}[nosep]
  \item \textbf{Rebuild}: first use or change of $\Delta t$ $\Rightarrow$ full
        setup (allocation + symbolic + numeric factorisation).
  \item \textbf{ValuesOnly}: same sparsity pattern, new values (the
        diagonal changes at every iteration in \dictkey{diagonalIion}) $\Rightarrow$ the
        values are refreshed reusing the pattern (\code{updateValues}).
  \item \textbf{Reuse}: matrix identical to the previous solve (Picard with a constant
        operator) $\Rightarrow$ the matrix is not touched, only solved.
\end{itemize}
The key \dictkey{petscReusePreconditioner} (\code{false} by default) additionally
allows, in \dictkey{diagonalIion} mode, reusing the preconditioner across
inner iterations when the diagonal change is small. In Picard this
optimisation is purely computational: the same matrix is solved with the same
factorisation, so the results are bit-for-bit identical
(Section~\ref{sec:exp-B}).

\subsection{Diffusion preconditioner for JFNK (A3)}
\label{sec:optim-A3}

The matrix-free JFNK Jacobian used to be solved without a preconditioner, so that
each Krylov iteration requires a full residual evaluation (and therefore
an ionic ODE integration). A \emph{right} preconditioner
$\mathbf{M}^{-1}$ is added with $\mathbf{M}=\Aimp$, the linear (diffusive, SPD) part
of the Jacobian. To avoid the catastrophic fill-in of an exact $LU$ factorisation
of the (wide-stencil) high-order operator, an \textbf{incomplete
$LU$ factorisation} (\code{Eigen::IncompleteLUT}, ILUT) is used, computed once and
reused; each Krylov step then costs a single cheap \emph{triangular solve}.
It is very effective here because, with small $\Delta t$, $\Aimp\approx
\Mmass/\Delta t$ is strongly diagonally dominant. The key
\dictkey{jfnkPreconditioner} takes the values \code{none} (default, legacy
path) or \code{diffusion}. In the PETSc branch \code{KSPSetPCSide(PC\_RIGHT)} is set
so that the KSP monitors the true residual $\norm{b - Ax}$; in the
in-house GMRES branch (Eigen) the preconditioner is applied in the Arnoldi step and in
the reconstruction of the correction. It does not change the converged solution, only the
convergence path (fewer matvecs, i.e.\ fewer ODE integrations).

\subsection{Parallel ionic ODE integration (C1)}
\label{sec:optim-C1}

The loop over the integration points of \code{TNNP::solveODE} is
\emph{embarrassingly parallel} (each point is an independent ODE). It is optionally
parallelised with OpenMP through the key \dictkey{parallelODE} (\code{false} by
default) read from the ionic-model dictionary. Since the base class
\code{ionicModel} shares a single \code{ODESolver} (with internal scratch for the
Runge--Kutta stages), a per-thread solver \emph{pool} is added
(\code{ensureOdeSolverPool}/\code{odeSolver(tid)}); the diagnostic index
\code{activeIntegrationPoint\_} becomes per-thread and writing the protection
log is serialised with \code{\#pragma omp critical}. It was necessary to add
\code{-fopenmp} to the compilation options of the \code{ionicModels} library.
The number of threads is controlled with the environment variable \code{OMP\_NUM\_THREADS};
\code{run\_convergence.py} exposes it through the constants \pyconst{PARALLEL\_ODE}
and \pyconst{ODE\_NUM\_THREADS}. The results are numerically identical to the
serial version (validated bit for bit).

\subsection{AMG preconditioners (C3)}
\label{sec:optim-C3}

The KSP path already supported passing PETSc options at run time, so
that the algebraic multigrid preconditioners (native \code{gamg}, \code{hypre}
BoomerAMG) are directly reachable. The aliases
\code{amg}$\to$\code{gamg} and \code{boomeramg}$\to$\code{hypre} are added to the
\dictkey{petscLinearPcType} mapping. For the SPD monodomain operator the
theoretically advisable combination on large 3D meshes is \code{cg}+\code{gamg}; however,
the experiments (Section~\ref{sec:exp-gamg}) show that for the time steps
used here the system is mass-dominated (well conditioned) and
AMG contributes nothing.

\paragraph{Summary of new keys.}
All with a default value equal to the previous behaviour.

\begin{center}
\begin{tabular}{@{}p{0.30\linewidth}p{0.17\linewidth}p{0.45\linewidth}@{}}
\toprule
Key & Default & Effect \\
\midrule
\dictkey{petscReusePreconditioner} & \code{false} & Reuses the preconditioner
    across iterations in \dictkey{diagonalIion}. \\
\dictkey{jfnkPreconditioner} & \code{none} & \code{diffusion}: preconditions the
    JFNK Jacobian with an ILUT of $\Aimp$. \\
\dictkey{parallelODE} & \code{false} & Integrates the ionic ODEs with OpenMP
    (ionic-model dict.). \\
\dictkey{petscLinearPcType} & \code{ilu} & Now accepts \code{gamg}/\code{hypre}
    (and the aliases \code{amg}/\code{boomeramg}). \\
\bottomrule
\end{tabular}
\end{center}

% ================================================================== %
\section{Performance experiments}
\label{sec:exp}
% ================================================================== %

The above optimisations were measured on the 2D Niederer-style benchmark
(\num{23726}-cell mesh unless stated otherwise) with the TNNP model and the RKF45
integrator, Crank--Nicolson scheme and consistent mass in the high-order cases.
Unless stated otherwise, the stimulus is the benchmark one (a
\SI{1.5}{\milli\meter} corner). The runs are short (few steps) and the measurements were
taken with \code{/usr/bin/time -v}; where relevant, \code{loopWall}
(time-loop time, from \code{solverPerformance.dat}) is reported to isolate the per-step
cost from the one-off setup. The machine has 20 cores.

\subsection{Phase B: identical and faster in Picard}
\label{sec:exp-B}

The same runs of the convergence sweep were compared against a binary
\emph{prior} to the changes (Picard, Crank--Nicolson, consistent mass, RKF45,
PETSc back-end \code{gmres}+\code{ilu}), taking the four cases present in
both folders.

\begin{center}
\begin{tabular}{@{}lccccc@{}}
\toprule
$\Delta t$ / $\Delta x$ & P8 [ms] (new=old) & nonlin.\ iter. & lin.\ iter. &
    wall new/old [s] & $\Delta$ \\
\midrule
$10^{-4}$ / \SI{0.25}{mm} & 62.9948077742 & 3510=3510 & 12748=12748 & 1488.7 / 1645.3 & $-9.5\%$ \\
$10^{-4}$ / \SI{0.5}{mm}  & 153.3156018900 & 6293=6293 & 18879=18879 & 684.0 / 752.0 & $-9.0\%$ \\
$10^{-5}$ / \SI{0.25}{mm} & 67.2950937434 & 6854=6854 & 13708=13708 & 2531.9 / 2786.9 & $-9.2\%$ \\
$10^{-5}$ / \SI{0.5}{mm}  & 157.1464148850 & 16136=16136 & 32272=32272 & 1472.6 / 1616.7 & $-8.9\%$ \\
\bottomrule
\end{tabular}
\end{center}

\noindent\textbf{Conclusion.} The P8 activation time agrees to ten
decimals ($\Delta=0$) and the number of nonlinear and linear iterations is
identical, while the wall time drops by $\sim 9\%$. Since CPU usage was
$99\%$ (serial), the saving comes entirely from not re-factorising the constant
matrix, not from parallelism.

\subsection{A3: diffusion preconditioner in JFNK}
\label{sec:exp-A3}

Same case (corner, 10 steps, JFNK, PETSc back-end), varying
\dictkey{jfnkPreconditioner}.

\begin{center}
\begin{tabular}{@{}lcc@{}}
\toprule
\dictkey{jfnkPreconditioner} & Krylov iter.\ (total) & wall [s] \\
\midrule
\code{none}      & 90 & 98.6 \\
\code{diffusion} & 24 & 59.0 \\
\bottomrule
\end{tabular}
\end{center}

\noindent\textbf{Conclusion.} The preconditioner reduces the matvecs from 90 to 24
($3.75\times$ fewer ODE integrations) and the wall time by $1.67\times$
($98.6\to59.0$~s), with the same solution (Vm agreeing within the Krylov
tolerance). Preconditioning is clearly faster than not doing it.

\subsection{C1: ODE parallelisation (does not help in this problem)}
\label{sec:exp-C1}

A/B test with \code{OMP\_NUM\_THREADS} fixed at 8, toggling \dictkey{parallelODE}
(the difference isolates the effect of C1, since the internal PETSc/BLAS threading is
the same in both). For the high-order and 3D cases \code{loopWall} is reported.

\begin{center}
\begin{tabular}{@{}lccc@{}}
\toprule
Configuration & \code{parallelODE=true} & \code{=false} & helps? \\
\midrule
Low-order Picard (2D)                  & 46.9  & 46.9  & no \\
JFNK + \code{diffusion} (2D)           & 57.06 & 57.58 & no \\
JFNK + \code{none} (2D)                & 28.76 & 28.96 & no \\
\code{gaussPointODE} + HO p3 (2D, loop)& 29.259 & 29.262 & no \\
3D, \num{17280} cells (loop)           & 2.861 & 2.816 & no \\
\bottomrule
\end{tabular}
\end{center}

\noindent The two scenarios where a benefit could be expected were specifically
tested: \dictkey{gaussPointODE} with high order p3 (several ODEs per cell,
$\sim\!16\times$ more ionic work) and 3D with many cells.
\textbf{Conclusion.} In every case the result is identical between
\code{true}/\code{false} (and the CPU $\%$ does not change), i.e.\ the parallel region
of the ODE contributes $\sim 0$. The reason is that the ionic ODE (RKF45 on TNNP, above
all at rest) is a \emph{small fraction} of the cost: what dominates is the LRE
reconstruction at the Gauss points, the averaging of $\Iion$ and the high-order
assembly/solve, which are serial operations that grow equally or more when the problem
is enlarged. C1 is functionally correct (bit-identical results) but does not speed up this
problem profile; it would only pay off with much stiffer models or extensive activation
fronts.

\subsection{Internal PETSc/BLAS threading (counterproductive in Picard)}
\label{sec:exp-threads}

Picard, \dictkey{parallelODE}=\code{false}, varying \code{OMP\_NUM\_THREADS}
(20 steps, corner).

\begin{center}
\begin{tabular}{@{}lccc@{}}
\toprule
\code{OMP\_NUM\_THREADS} & wall [s] & CPU & User [s] \\
\midrule
1  & 45.7 & $98\%$  & 44.9 \\
20 & 47.3 & $197\%$ & 68.8 \\
\bottomrule
\end{tabular}
\end{center}

\noindent\textbf{Conclusion.} With more threads, internal PETSc/BLAS/Eigen
parallelism appears ($\sim 2\times$ CPU), but the wall time does not drop (it rises
slightly): it is \emph{memory-bound} and only wastes CPU. Picard should be run
with \code{OMP\_NUM\_THREADS=1}.

\subsection{Linear solver: AMG vs.\ ILU (the solve is not the bottleneck)}
\label{sec:exp-gamg}

Comparison of \dictkey{petscLinearKspType}/\dictkey{petscLinearPcType} with
\code{OMP\_NUM\_THREADS=1}.

\begin{center}
\begin{tabular}{@{}llcc@{}}
\toprule
Case & Linear solver & \code{loopWall} [s] & lin.\ iter.\ (total) \\
\midrule
3D, \num{17280}, $\Delta t{=}10^{-6}$, low order
   & \code{gmres}+\code{ilu}  & 5.27 & 16 \\
   & \code{cg}+\code{gamg}    & 5.72 & 24 \\
\midrule
2D, \num{23726}, $\Delta t{=}10^{-4}$, HO consistent
   & \code{gmres}+\code{ilu}  & 70.21 & 40 \\
   & \code{gmres}+\code{gamg} & 70.66 & 72 \\
   & \code{cg}+\code{gamg}    & 70.47 & 72 \\
\bottomrule
\end{tabular}
\end{center}

\noindent\textbf{Conclusion.} AMG is slightly slower and uses \emph{more}
iterations: with small $\Delta t$ the system is mass-dominated
($\Aimp\approx\Mmass/\Delta t$), well conditioned and reaction-like, where ILU already
is an excellent local preconditioner and the multigrid hierarchy is
counterproductive. More importantly still: the \code{loopWall} is practically the same
($\sim 70$~s) whatever the linear solver (40 vs.\ 72 iterations), which
shows that \textbf{the linear solve is not the bottleneck}. Together with the
previous experiments, the dominant cost in the high-order Picard regime is the
high-order machinery itself (LRE reconstruction and Gauss-point handling),
not the ODE nor the linear solver.

\subsection{Synthesis and recommendations}
\label{sec:exp-sintesis}

\begin{itemize}[nosep]
  \item Keep \dictkey{parallelODE}=\code{false} and \code{OMP\_NUM\_THREADS=1}:
        neither C1 nor the internal threading speeds up this problem.
  \item Keep \dictkey{petscLinearPcType}=\code{ilu} (do not use AMG): the system
        is well conditioned and the linear solve is not the bottleneck.
  \item The real gains, without parallelism, are Phase~B
        (factorisation reuse, $\sim 9\%$ in Picard, identical
        results) and A3 (\dictkey{jfnkPreconditioner}=\code{diffusion},
        $1.67\times$ in JFNK).
  \item The natural target to speed up the high-order regime would be to parallelise
        or make cheaper the LRE reconstruction and the Gauss-point handling (not
        done here).
\end{itemize}

% ================================================================== %
\begin{thebibliography}{9}
% ================================================================== %

\bibitem{ttp2004} K.~H.~W.~J. ten Tusscher, D. Noble, P.~J. Noble, A.~V.
    Panfilov, \emph{A model for human ventricular tissue},
    AJP -- Heart and Circulatory Physiology
    \textbf{286} (2004) H1573--H1589.

\bibitem{niederer2012} S.~A. Niederer et al.,
    \emph{Verification of cardiac tissue electrophysiology simulators
    using an N-version benchmark}, Prog.\ Biophys.\ Mol.\ Biol.
    \textbf{104} (2012) 22--48.

\bibitem{knoll2004} D.~A. Knoll, D.~E. Keyes,
    \emph{Jacobian-Free Newton-Krylov methods: a survey of approaches
    and applications}, J.\ Comp.\ Phys.\ \textbf{193}
    (2004) 357--397.

\bibitem{noceWright} J. Nocedal, S.~J. Wright,
    \emph{Numerical Optimization}, 2nd edition, Springer, 2006.

\bibitem{saadbook} Y. Saad,
    \emph{Iterative Methods for Sparse Linear Systems}, 2nd edition,
    SIAM, 2003.

\bibitem{hairer2} E. Hairer, G. Wanner,
    \emph{Solving Ordinary Differential Equations II}, Springer, 1996.

\end{thebibliography}

\end{document}
